% Scoring ledger: Q1/Q2 C1=2+2=4, C2=1+2+1=4; Q3 C1=3+3=6, C2=2+2+2=6.
% Totals: C1=14, C2=14; overall=28. Construction validity is within Q3 marks.
% Governing source: G11_T2_L1_C1C2_exam.tex (current exam).
% Reference activities: ../2-class_activities/G11_T2_L1_C1_classroom_activity.tex
% and G11_T2_L1_C2_classroom_activity.tex. No archived exam is governing.
% Build from this directory; TEXINPUTS must include ../../../preamble/.
\makeatletter
\def\input@path{{../../../preamble/}}
\makeatother
\input{preamble_exams.tex}
\usepackage{needspace}
\input{highlights/highlights.tex}
% \GTenHighlightsfalse % Uncomment to disable instructional highlights.
\input{highlights/specific_highlights.tex}
\newcommand{\CtwoSignificance}[1]{\CtwoConfidence{#1}}
\definecolor{CtwoErrorColor}{HTML}{28AFA2}
\newcommand{\CtwoError}[1]{\HighlightStroke{CtwoErrorColor}{#1}}
\newcommand{\CtwoCI}[1]{\CtwoEndpoint{#1}}
\newcommand{\Pop}[1]{\HighlightStroke{G11T2C1Population}{#1}}
\newcommand{\Sam}[1]{\HighlightStroke{G11T2C1Sample}{#1}}
\newcommand{\SSD}[1]{\HighlightStroke{G11T2C1SampleSD}{#1}}
\newcommand{\PP}[1]{\HighlightStroke{G11T2C1PopulationProportion}{#1}}
\newcommand{\SP}[1]{\HighlightStroke{G11T2C1SampleProportion}{#1}}

\hypersetup{pdftitle={Grade 11 Term 2 C1 and C2 Exam Solutions},pdfauthor={Global Economics}}
\raggedbottom
\linespread{1.04}
\setlength{\parskip}{4pt}
\newcolumntype{K}{>{\raggedright\arraybackslash}X}
\newcommand{\keytable}[2]{\par\noindent\begingroup\small\setlength{\tabcolsep}{4pt}\renewcommand{\arraystretch}{1.14}%
\begin{tabularx}{\linewidth}{#1}\toprule #2\bottomrule\end{tabularx}\par\endgroup}
\newcommand{\examtable}[2]{\par\noindent\begin{minipage}{\linewidth}\small
\renewcommand{\arraystretch}{1.08}\setlength{\tabcolsep}{3pt}
\begin{tabularx}{\linewidth}{#1}\toprule #2\bottomrule\end{tabularx}
\end{minipage}\par}
\newcommand{\taskprompt}[3]{\Needspace{5\baselineskip}\par\textbf{Original task: #1}\score{#2}\par #3\par}
\newenvironment{probleminfoexcerpt}{%
  \begin{list}{}{\setlength{\leftmargin}{1.5em}\setlength{\rightmargin}{1.5em}%
    \setlength{\topsep}{3pt}\setlength{\partopsep}{0pt}\setlength{\parsep}{0pt}}%
  \item\relax\small\raggedright}{\end{list}}
\newcommand{\worked}[2]{\subsection*{Worked solution: #1}\label{#2}}
\newcommand{\cOneA}{Identify the population and both samples: place, period and sizes. Pair and interpret population and sample means, standard deviations and proportions; state which are known. Calculate both sample proportions from their condition counts.}
\newcommand{\cOneB}{Calculate second-minus-initial changes in mean, sample standard deviation and proportion (percentage points). Explain which population characteristics remain fixed and why sample differences do not establish parameter changes.}
\newcommand{\cTwoA}{Complete the C2 information table for the initial study: initially unknown and later established mean, known spread, sample statistics and procedure settings. Calculate SE and give confidence, significance, each tail area and critical value at every requested level.}
\newcommand{\cTwoB}{For each requested confidence level, tabulate the initial point estimate, margin of error, separate endpoints, confidence interval and full width using the known population standard deviation.}
\newcommand{\cTwoC}{For subsequent studies only, calculate covering counts, observed coverage and absolute differences from nominal confidence (percentage points), including endpoint equality. Explain long-run coverage and finite-study variation; distinguish coverage from sample means inside the initial interval and from probability assigned to a fixed mean after calculation.}
\title{Grade 11 -- Term 2\\C1 and C2 Integrated Exam --- Solutions}
\author{Global Economics}
\date{2026--2027}
\begin{document}
\maketitle
\setcounter{tocdepth}{1}
\phantomsection\label{toc}\tableofcontents
\section*{Learning objective}
Distinguish fixed population parameters from variable sample statistics in economic and financial contexts; calculate known-population-standard-deviation confidence intervals and interpret their long-run coverage.
\section*{Criteria assessment}
Complete all three problems. Problems 1 and 2 each award 4 points to C1 and 4 to C2 (8 overall). Problem 3 awards 6 points to each criterion (12 overall). C1 has 14 available points, C2 has 14, and the exam has 28.

\textbf{Each criterion requires at least 4 out of 14 points from this exam to pass independently. Points cannot transfer between criteria.} An overall total alone does not establish a pass in both criteria.

Show reasoning, formulas, substitutions, units and contextual interpretations. Use the reference table formats. Give exact probabilities where applicable. Retain all endpoint and width digits generated by the supplied rounded critical values. Label any graphs; monetary axes must use actual monetary spacing. All scenario data are hypothetical.
\section*{Evaluated criteria}
\begin{description}
\item[C1] Distinguishes population parameters from sample statistics by interpreting and comparing population and sample means, standard deviations and proportions in economic and financial contexts.
\item[C2] Compares point and interval estimation and interprets confidence through long-run coverage, distinguishing this from probability assigned to a particular calculated interval containing the parameter.
\end{description}
\clearpage
\textbf{Use of this key.} Each original question is printed with its data; every task appears before its worked solution. Links lead to solutions; return arrows lead to the question and contents. All detailed marking breakdowns below are proposed partial-credit guidance under the prescribed task marks. Hypothetical data are used throughout.
\section*{General tasks for every problem}
Complete \textbf{all five tasks for each of the three problems}, using the supplied scenario and data in Problems 1 and 2 and your own constructed scenario and data in Problem 3. Use the confidence level or levels specified in each scenario. The shared table after these five tasks gives the marks for each problem.
\begin{enumerate}[label=(\alph*),itemsep=1pt]
\item \textbf{C1 Task 1}. \cOneA
\item \textbf{C1 Task 2}. \cOneB
\item \textbf{C2 Task 1}. \cTwoA
\item \textbf{C2 Task 2}. \cTwoB
\item \textbf{C2 Task 3}. \cTwoC
\end{enumerate}

\examtable{Krrr}{
\textbf{Task}&\textbf{Problem 1}&\textbf{Problem 2}&\textbf{Problem 3}\\\midrule
C1 Task 1&2&2&3\\
C1 Task 2&2&2&3\\
C2 Task 1&1&1&2\\
C2 Task 2&2&2&2\\
C2 Task 3&1&1&2\\
\midrule
\textbf{C1 subtotal}&\textbf{4}&\textbf{4}&\textbf{6}\\
\textbf{C2 subtotal}&\textbf{4}&\textbf{4}&\textbf{6}\\
\textbf{Problem total}&\textbf{8}&\textbf{8}&\textbf{12}\\}

\clearpage
\section*{Task reference material}
\textbf{C1 pairings.} Population $\mathcal P$, size $N$; sample $X$, size $n$; population/sample mean $(\mu,\bar x)$; standard deviation $(\sigma,s)$; proportion $(p,\hat p)$. A condition count is $k$, so $\hat p=k/n$. A sample-proportion change is $100(\hat p_B-\hat p_A)$ percentage points.

\textbf{C2 information table.} Use contiguous groups in this order: Population parameters ($\mu,\mu_{\mathrm{known}},\sigma$); Sample statistics ($\bar x,n,\mathrm{SE}$); Procedure settings ($C=1-\alpha,\alpha,\alpha/2,z_{1-\alpha/2}$). Columns are Symbol, Concept, Type and Meaning in context. As in the classroom reference, the derived SE is placed in the sample-statistics group; here it uses known $\sigma$ and fixed $n$, so it is not an estimate using $s$. The later-known notation refers to the same fixed mean, not a second parameter.

\textbf{Two-sided standard normal reference.} Use the following numerical critical values for all calculations. They are supplied rounded values; retain all digits generated by them when calculating endpoints and widths.
\par\noindent\begin{minipage}{\linewidth}
\begin{minipage}[t]{.49\linewidth}\vspace{0pt}
\keytable{@{}KK@{}}{
\textbf{Confidence level}\newline $1-\alpha$&\textbf{Positive critical value}\newline $z_{1-\alpha/2}$\\\midrule
80\%&1.282\\
90\%&1.645\\
95\%&1.960\\
}
\end{minipage}\hfill
\begin{minipage}[t]{.49\linewidth}\vspace{0pt}
\keytable{@{}KK@{}}{
\textbf{Confidence level}\newline $1-\alpha$&\textbf{Positive critical value}\newline $z_{1-\alpha/2}$\\\midrule
98\%&2.326\\
99\%&2.576\\
99.9\%&3.291\\
}
\end{minipage}
\end{minipage}\par
\[
\begin{aligned}
\alpha&=1-C,& \text{each tail}&=\alpha/2,&\mathrm{SE}&=\frac{\sigma}{\sqrt n},\\
\mathrm{MoE}&=z_{1-\alpha/2}\mathrm{SE},&L&=\bar x-\mathrm{MoE},&U&=\bar x+\mathrm{MoE},\\
\mathrm{CI}&=[L,U],&W&=U-L=2\mathrm{MoE}.&&
\end{aligned}
\]
For each confidence level separately, a subsequent interval covers when
$L_i\leq\mu_{\mathrm{known}}\leq U_i$. Equality at an endpoint counts as coverage.
\[
\mathrm{Cov}(C)=\frac{N_{\mathrm{cover}}(C)}{N(C)}\,100\%,\qquad
E_{\mathrm{abs}}(C)=|\text{observed percentage}-100C|\text{ percentage points}.
\]
The interval table uses Symbol, Concept, Calculation and Result and meaning in context, with separate endpoint rows. Proportions and confidence settings are unitless; a mean, standard deviation, standard error, margin, endpoint and width use the measurement's units.

\textbf{Common procedure assumptions.} The scenarios specify independent random studies, negligible within-study sampling fractions and sufficiently regular numerical populations for a normal approximation to the sampling distribution of the mean. The finite-population correction is negligible and omitted. Thus nominal normal coverage is an approximation, additionally subject to supplied critical-value rounding. Sample standard deviations describe samples; they do not replace known population standard deviations in this procedure.
\secInicio

\clearpage
\section*{Instructional highlight guide}\phantomsection\label{key:legend}
\addcontentsline{toc}{section}{Instructional highlight guide}
The C1 and C2 classroom colors connect concept names, notation and values. Labels and all content remain when highlights are disabled. Detailed contextual meanings appear in each worked information and calculation table.
\keytable{p{.24\linewidth}p{.28\linewidth}K}{
\textbf{Symbol or notation}&\textbf{Concept}&\textbf{Brief meaning}\\\midrule
$\Pop{\mathcal P,N}$&\Pop{Population and size}&All eligible units in the fixed scope.\\
$\Sam{X_A,X_B,n,k}$&\Sam{Samples, size and count}&Selected units and condition counts.\\
$\CtwoMean{\mu,\mu_{\mathrm{known}}}$&\CtwoMean{Population mean}&Same fixed mean; later established value.\\
$\CtwoPoint{\bar x_A,\bar x_B}$&\CtwoPoint{Sample mean}&Average within each sample; point estimate.\\
$\CtwoSD{\sigma}$&\CtwoSD{Population SD}&Known population spread.\\
$\SSD{s_A,s_B}$&\SSD{Sample SD}&Spread within each sample.\\
$\PP{p}$&\PP{Population proportion}&Population condition share.\\
$\SP{\hat p_A,\hat p_B}$&\SP{Sample proportion}&Sample condition shares.\\
$\CtwoSE{\mathrm{SE}}$&\CtwoSE{Standard error}&Spread of sample means.\\
$\CtwoConfidence{C,\alpha,\alpha/2}$&\CtwoConfidence{Confidence and tail settings}&Central coverage, significance and tail areas.\\
$\CtwoCritical{z_{1-\alpha/2}}$&\CtwoCritical{Critical value}&Supplied standard normal cutoff.\\
$\CtwoMoE{\mathrm{MoE}}$&\CtwoMoE{Margin of error}&Half-width of the interval.\\
$\CtwoEndpoint{L,U,\mathrm{CI}}$&\CtwoEndpoint{Endpoints and interval}&Bounds and range for a mean.\\
$\CtwoWidth{W}$&\CtwoWidth{Full width}&Distance between endpoints.\\
$\CtwoCoverage{N_{\mathrm{cover}},\mathrm{Cov}}$&\CtwoCoverage{Observed coverage}&Covering count and percentage.\\
$\CtwoError{E_{\mathrm{abs}}}$&\CtwoError{Absolute difference}&Distance from nominal confidence in percentage points.\\}
\secInicio

\clearpage

\sectionScored{Retail electricity costs}{C1-\sml{4}\quad C2-\sml{4}}{sec:p1}

\textbf{Original shared description}\score{8 overall}

An energy-cost review concerns all 200,000 registered retail businesses operating in Medellin throughout October 2026. Each business contributes its total October electricity bill, measured in thousand COP. A separate yes-or-no condition records whether the business used electronic invoicing in October.

The initial independent random sample contains 100 businesses. Its mean bill is 240 thousand COP, its sample standard deviation is 58 thousand COP, and 40 businesses used electronic invoicing. A second independent random sample from exactly the same business population and October period also contains 100 businesses. Its mean bill is 252 thousand COP, its sample standard deviation is 66 thousand COP, and 55 businesses used electronic invoicing.

A certified spread-only population report makes the population standard deviation of October bills known as 60 thousand COP; it does not disclose the population mean. The two sample standard deviations remain separate descriptive quantities. The initial study estimates the initially unknown mean at 95\% confidence.

Assume the bill distribution is sufficiently regular, without dominant outliers, for a normal approximation to sample means of this size. Each study draws a simple random sample without replacement from the full fixed register. Samples are selected independently across studies; overlap between different samples is permitted. Each study selects only 0.05\% of the population, so dependence within a draw and the finite-population correction are negligible.

Ten subsequent independent studies each use 100 businesses, the same known population standard deviation and the same 95\% normal confidence procedure. These are additional studies, separate from both samples above. A later comprehensive investigation of the entire October population establishes its mean bill as 245 thousand COP. Repeated samples do not establish this value. The following subsequent intervals allow retrospective assessment; all endpoints are in thousand COP.

\par\noindent\begin{minipage}{\linewidth}
\begin{minipage}[t]{.49\linewidth}
\keytable{@{}>{\hsize=1.3\hsize}K@{\hspace{2pt}}>{\hsize=.85\hsize}K@{\hspace{2pt}}>{\hsize=.85\hsize}K@{}}{
\textbf{Subsequent study}&\textbf{Lower endpoint}&\textbf{Upper endpoint}\\\midrule
1&224.24&247.76\\
2&232.24&255.76\\
3&238.24&261.76\\
4&246.24&269.76\\
5&218.24&241.76\\
}
\end{minipage}\hfill
\begin{minipage}[t]{.49\linewidth}
\keytable{@{}>{\hsize=1.3\hsize}K@{\hspace{2pt}}>{\hsize=.85\hsize}K@{\hspace{2pt}}>{\hsize=.85\hsize}K@{}}{
\textbf{Subsequent study}&\textbf{Lower endpoint}&\textbf{Upper endpoint}\\\midrule
6&235.24&258.76\\
7&230.24&253.76\\
8&242.24&265.76\\
9&221.24&244.76\\
10&234.24&257.76\\
}
\end{minipage}
\end{minipage}\par


Complete all five \textbf{General tasks for every problem} above; see the shared table for task marks.

\clearpage
\textbf{Linked tasks.} Each full task is reproduced immediately before its solution.

\begin{enumerate}[label=(\alph*),itemsep=6pt]
\item \hyperref[sol:p1c1a]{\textbf{C1 Task 1}}. \cOneA\ \sscore{\sml{2}}
\item \hyperref[sol:p1c1b]{\textbf{C1 Task 2}}. \cOneB\ \sscore{\sml{2}}
\item \hyperref[sol:p1c2a]{\textbf{C2 Task 1}}. \cTwoA\ \sscore{\sml{1}}
\item \hyperref[sol:p1c2b]{\textbf{C2 Task 2}}. \cTwoB\ \sscore{\sml{2}}
\item \hyperref[sol:p1c2c]{\textbf{C2 Task 3}}. \cTwoC\ \sscore{\sml{1}}
\end{enumerate}
\secInicio[][sec:p1]

\clearpage

\taskprompt{C1 Task 1}{2}{\cOneA}

\worked{C1 Task 1 \score{2}}{sol:p1c1a}

\textbf{Relevant problem information}

\begin{probleminfoexcerpt}
An energy-cost review concerns all \Pop{200,000} \Pop{registered} \Pop{retail} \Pop{businesses} operating in Medellin throughout October 2026. Each business contributes its total October electricity bill, measured in thousand COP. A separate yes-or-no condition records whether the business used electronic invoicing in October.
\end{probleminfoexcerpt}

\begin{probleminfoexcerpt}
The \Sam{initial} \Sam{independent} \Sam{random} \Sam{sample} contains \Sam{100} businesses. Its \CtwoPoint{mean} \CtwoPoint{bill} is \CtwoPoint{240} \CtwoPoint{thousand} \CtwoPoint{COP}, its \SSD{sample} \SSD{standard} \SSD{deviation} is \SSD{58} \SSD{thousand} \SSD{COP}, and \Sam{40} businesses used electronic invoicing. A \Sam{second} \Sam{independent} \Sam{random} \Sam{sample} from exactly the same business population and October period also contains \Sam{100} businesses. Its \CtwoPoint{mean} \CtwoPoint{bill} is \CtwoPoint{252} \CtwoPoint{thousand} \CtwoPoint{COP}, its \SSD{sample} \SSD{standard} \SSD{deviation} is \SSD{66} \SSD{thousand} \SSD{COP}, and \Sam{55} businesses used electronic invoicing.
\end{probleminfoexcerpt}

\begin{probleminfoexcerpt}
A certified spread-only population report makes the \CtwoSD{population} \CtwoSD{standard} \CtwoSD{deviation} of \CtwoSD{October bills} known as \CtwoSD{60} \CtwoSD{thousand} \CtwoSD{COP}; it does not disclose the \CtwoMean{population} \CtwoMean{mean October} \CtwoMean{electricity bill}.
\end{probleminfoexcerpt}

\textbf{Later investigation (retrospective information).}

\begin{probleminfoexcerpt}
A later comprehensive investigation of the entire October population establishes its \CtwoMean{mean October} \CtwoMean{electricity bill} as \CtwoMean{245} \CtwoMean{thousand} \CtwoMean{COP}. Repeated samples do not establish this value.
\end{probleminfoexcerpt}

They are separate independent selections, not populations at two different dates.

\textbf{Paired population and sample measures.} All monetary measures below use thousand COP. Sample proportions describe the condition defined in the population table.
\keytable{p{.12\linewidth}p{.24\linewidth}K}{
\textbf{Symbol}&\textbf{Concept}&\textbf{Meaning in context}\\\midrule
$\Pop{\mathcal P,\ N}$&Population and size&\Pop{200,000 eligible businesses} in the scope stated above.\\
$\CtwoMean{\mu}$&\CtwoMean{Population mean}&\CtwoMean{mean October} \CtwoMean{electricity bill}; initially unknown, later $\CtwoMean{245}$ thousand COP.\\
$\CtwoSD{\sigma}$&\CtwoSD{Population SD}&$\CtwoSD{60}$ thousand COP; known \CtwoSD{population standard} \CtwoSD{deviation}, the \CtwoSD{spread of October} \CtwoSD{electricity bills}.\\
$\PP{p}$&\PP{Population proportion}&\PP{proportion of all} \PP{eligible businesses} \PP{that used electronic} \PP{invoicing for their} \PP{October electricity} \PP{bills}; unknown.\\}
\keytable{p{.16\linewidth}p{.23\linewidth}KK}{
\textbf{Symbol}&\textbf{Concept}&\textbf{Initial sample A}&\textbf{Second sample B}\\\midrule
$\Sam{X_A,X_B;\ n_A,n_B}$&Samples and sizes&$\Sam{100}$ randomly selected businesses&$\Sam{100}$ independently selected businesses\\
$\CtwoPoint{\bar x_A,\bar x_B}$&\CtwoPoint{Sample mean}&$\CtwoPoint{240}$ thousand COP&$\CtwoPoint{252}$ thousand COP\\
$\SSD{s_A,s_B}$&\SSD{Sample SD}&$\SSD{58}$ thousand COP&$\SSD{66}$ thousand COP\\
$\Sam{k_A,k_B}$&Condition count&$\Sam{40}$ businesses&$\Sam{55}$ businesses\\
$\SP{\hat p_A,\hat p_B}$&\SP{Sample proportion}&$\SP{\hat p_A}=\Sam{40}/\Sam{100}=\SP{0.4=40\%}$&$\SP{\hat p_B}=\Sam{55}/\Sam{100}=\SP{0.55=55\%}$\\}



The condition counts are $\Sam{k_A}=\Sam{40}$ and $\Sam{k_B}=\Sam{55}$. The proportions are unitless shares of each sample meeting the stated condition. The known \CtwoSD{population SD} does not make either \SSD{sample SD} equal to it. Initially $\CtwoMean{\mu}$ is unknown; the later investigation establishes that same parameter as $\CtwoMean{245}$ thousand COP. The supplied information never establishes $\PP{p}$: sample proportions are estimates, and the investigation reports only the population mean.

\secInicio[][sec:p1]

\clearpage

\taskprompt{C1 Task 2}{2}{\cOneB}

\worked{C1 Task 2 \score{2}}{sol:p1c1b}

\textbf{Relevant problem information}

\begin{probleminfoexcerpt}
An energy-cost review concerns all \Pop{200,000} \Pop{registered} \Pop{retail} \Pop{businesses} operating in Medellin throughout October 2026. Each business contributes its total October electricity bill, measured in thousand COP. A separate yes-or-no condition records whether the business used electronic invoicing in October.
\end{probleminfoexcerpt}

\begin{probleminfoexcerpt}
The \Sam{initial} \Sam{independent} \Sam{random} \Sam{sample} contains \Sam{100} businesses. Its \CtwoPoint{mean} \CtwoPoint{bill} is \CtwoPoint{240} \CtwoPoint{thousand} \CtwoPoint{COP}, its \SSD{sample} \SSD{standard} \SSD{deviation} is \SSD{58} \SSD{thousand} \SSD{COP}, and \Sam{40} businesses used electronic invoicing. A \Sam{second} \Sam{independent} \Sam{random} \Sam{sample} from exactly the same business population and October period also contains \Sam{100} businesses. Its \CtwoPoint{mean} \CtwoPoint{bill} is \CtwoPoint{252} \CtwoPoint{thousand} \CtwoPoint{COP}, its \SSD{sample} \SSD{standard} \SSD{deviation} is \SSD{66} \SSD{thousand} \SSD{COP}, and \Sam{55} businesses used electronic invoicing.
\end{probleminfoexcerpt}

\keytable{KKK}{
\textbf{Second minus initial}&\textbf{Calculation}&\textbf{Interpretation}\\\midrule
\CtwoPoint{Mean}&$\CtwoPoint{252-240=+12}$ thousand COP&Second sample mean is higher.\\
\SSD{Sample SD}&$\SSD{66-58=+8}$ thousand COP&Second sample has greater numerical spread.\\
\SP{Proportion}&$\SP{55\%-40\%=15}$ percentage points&Second sample has a larger condition share.\\}


Both sample sizes remain 100; the selected units and descriptive statistics differ. Population membership, size, place, period, measurement and binary-condition definitions remain fixed. Therefore $\CtwoMean{\mu}$, $\CtwoSD{\sigma}$ and $p$ remain fixed, whether known or unknown. Random sampling variation can produce these differences; they do not establish a temporal change in population parameters. Learning the value of $\CtwoMean{\mu}$ later changes our knowledge, not the population mean itself.

\secInicio[][sec:p1]

\clearpage

\taskprompt{C2 Task 1}{1}{\cTwoA}

\worked{C2 Task 1 \score{1}}{sol:p1c2a}

\textbf{Relevant problem information}

\begin{probleminfoexcerpt}
An energy-cost review concerns all \Pop{200,000} \Pop{registered} \Pop{retail} \Pop{businesses} operating in Medellin throughout October 2026. Each business contributes its total October electricity bill, measured in thousand COP.
\end{probleminfoexcerpt}

\begin{probleminfoexcerpt}
The \CtwoSize{initial} \CtwoSize{independent} \CtwoSize{random} \CtwoSize{sample} contains \CtwoSize{100} businesses. Its \CtwoPoint{mean} \CtwoPoint{bill} is \CtwoPoint{240} \CtwoPoint{thousand} \CtwoPoint{COP}, its \SSD{sample} \SSD{standard} \SSD{deviation} is \SSD{58} \SSD{thousand} \SSD{COP}, and \Sam{40} businesses used electronic invoicing.
\end{probleminfoexcerpt}

\begin{probleminfoexcerpt}
A certified spread-only population report makes the \CtwoSD{population} \CtwoSD{standard} \CtwoSD{deviation} of \CtwoSD{October bills} known as \CtwoSD{60} \CtwoSD{thousand} \CtwoSD{COP}; it does not disclose the \CtwoMean{population} \CtwoMean{mean October} \CtwoMean{electricity bill}. The initial study estimates the initially unknown \CtwoMean{mean October} \CtwoMean{electricity bill} at \CtwoConfidence{95\%} \CtwoConfidence{confidence}.
\end{probleminfoexcerpt}

\begin{probleminfoexcerpt}
Assume the bill distribution is sufficiently regular, without dominant outliers, for a normal approximation to sample means of this size. Each study draws a simple random sample without replacement from the full fixed register. Samples are selected independently across studies; overlap between different samples is permitted. Each study selects only 0.05\% of the population, so dependence within a draw and the finite-population correction are negligible.
\end{probleminfoexcerpt}

\textbf{Later investigation (unavailable for initial estimation).}

\begin{probleminfoexcerpt}
A later comprehensive investigation of the entire October population establishes its \CtwoMean{mean October} \CtwoMean{electricity bill} as \CtwoMean{245} \CtwoMean{thousand} \CtwoMean{COP}. Repeated samples do not establish this value.
\end{probleminfoexcerpt}

\begingroup\raggedright The unknown $\PP{p}$ is the \PP{proportion of all} \PP{eligible businesses} \PP{that used electronic} \PP{invoicing for their} \PP{October electricity} \PP{bills}. Its denominator is all 200,000 eligible businesses defined above.\par\endgroup

\keytable{p{.12\linewidth}p{.24\linewidth}K}{
\textbf{Symbol}&\textbf{Concept}&\textbf{Meaning in context}\\\midrule
\multicolumn{3}{l}{\textbf{Population parameters}}\\
$\CtwoMean{\mu}$&\CtwoMean{Population mean}&Initially unknown \CtwoMean{mean October} \CtwoMean{electricity bill}.\\
$\CtwoMean{\mu_{\mathrm{known}}}$&\CtwoMean{Later established} \CtwoMean{mean}&$\CtwoMean{245}$ thousand COP; the same fixed \CtwoMean{mean October} \CtwoMean{electricity bill}, established by the population investigation; unavailable for initial estimation.\\
$\CtwoSD{\sigma}$&\CtwoSD{Population SD}&$\CtwoSD{60}$ thousand COP; known \CtwoSD{population standard} \CtwoSD{deviation}, the \CtwoSD{spread of October} \CtwoSD{electricity bills}.\\\midrule
\multicolumn{3}{l}{\textbf{Sample statistics}}\\
$\CtwoPoint{\bar x}$&\CtwoPoint{Initial sample mean}&$\CtwoPoint{240}$ thousand COP; point estimate.\\
$\CtwoSize{n}$&\CtwoSize{Initial sample size}&$\CtwoSize{100}$ businesses.\\
$\CtwoSE{\mathrm{SE}}$&\CtwoSE{Standard error}&$\CtwoSE{\mathrm{SE}}=\frac{\CtwoSD{\sigma}}{\sqrt{\CtwoSize{n}}}=\frac{\CtwoSD{60}}{\sqrt{\CtwoSize{100}}}=\CtwoSE{6}$ thousand COP; spread of sample means, using known $\CtwoSD{\sigma}$ and the initial sample size $\CtwoSize{n}$.\\\midrule
\multicolumn{3}{l}{\textbf{Procedure settings}}\\
$\CtwoConfidence{C}$&\CtwoConfidence{Confidence}&$\CtwoConfidence{95\%}$: $\CtwoConfidence{0.95}$.\\
$\CtwoSignificance{\alpha}$&\CtwoSignificance{Significance}&$\CtwoConfidence{95\%}$: $1-\CtwoConfidence{0.95}=\CtwoSignificance{0.05}$.\\
$\CtwoSignificance{\alpha/2}$&\CtwoSignificance{Each tail area}&$\CtwoConfidence{95\%}$: $\CtwoSignificance{0.05}/2=\CtwoSignificance{0.025}$.\\
$\CtwoCritical{z_{1-\alpha/2}}$&\CtwoCritical{Positive} \CtwoCritical{critical value}&$\CtwoConfidence{95\%}$: $\CtwoCritical{1.960}$.\\}



Use the initial sample only for estimation. The comparison sample describes sample variation; the later mean is reserved for retrospective assessment. Known $\CtwoSD{\sigma}=\CtwoSD{60}$ is used, rather than either $\SSD{s_A}=\SSD{58}$ or $\SSD{s_B}=\SSD{66}$. The normal approximation is supported by random sampling, the stated regularity assumption and negligible sample fraction.

\secInicio[][sec:p1]

\clearpage

\taskprompt{C2 Task 2}{2}{\cTwoB}

\worked{C2 Task 2 \score{2}}{sol:p1c2b}

\textbf{Relevant problem information}

\begin{probleminfoexcerpt}
An energy-cost review concerns all \Pop{200,000} \Pop{registered} \Pop{retail} \Pop{businesses} operating in Medellin throughout October 2026. Each business contributes its total October electricity bill, measured in thousand COP.
\end{probleminfoexcerpt}

\begin{probleminfoexcerpt}
The \CtwoSize{initial} \CtwoSize{independent} \CtwoSize{random} \CtwoSize{sample} contains \CtwoSize{100} businesses. Its \CtwoPoint{mean} \CtwoPoint{bill} is \CtwoPoint{240} \CtwoPoint{thousand} \CtwoPoint{COP}, its \SSD{sample} \SSD{standard} \SSD{deviation} is \SSD{58} \SSD{thousand} \SSD{COP}, and \Sam{40} businesses used electronic invoicing.
\end{probleminfoexcerpt}

\begin{probleminfoexcerpt}
A certified spread-only population report makes the \CtwoSD{population} \CtwoSD{standard} \CtwoSD{deviation} of October bills known as \CtwoSD{60} \CtwoSD{thousand} \CtwoSD{COP}; it does not disclose the \CtwoMean{population} \CtwoMean{mean}. The initial study estimates the initially unknown mean at \CtwoConfidence{95\%} \CtwoConfidence{confidence}.
\end{probleminfoexcerpt}

\begin{probleminfoexcerpt}
Assume the bill distribution is sufficiently regular, without dominant outliers, for a normal approximation to sample means of this size. Each study draws a simple random sample without replacement from the full fixed register. Samples are selected independently across studies; overlap between different samples is permitted. Each study selects only 0.05\% of the population, so dependence within a draw and the finite-population correction are negligible.
\end{probleminfoexcerpt}

\textbf{Later investigation (unavailable for initial estimation).}

\begin{probleminfoexcerpt}
A later comprehensive investigation of the entire October population establishes its \CtwoMean{mean} \CtwoMean{bill} as \CtwoMean{245} \CtwoMean{thousand} \CtwoMean{COP}. Repeated samples do not establish this value.
\end{probleminfoexcerpt}

The shared \CtwoPoint{point estimate} is $\CtwoPoint{\bar x}=\CtwoPoint{240}$ thousand COP, the average October electricity bill in the initial sample; it estimates the population mean with a single value.

\textbf{\CtwoConfidence{95\%} confidence}

\keytable{p{.10\linewidth}p{.17\linewidth}p{.32\linewidth}K}{
\textbf{Symbol}&\textbf{Concept}&\textbf{Calculation}&\textbf{Result and meaning in context}\\\midrule
$\CtwoPoint{\bar x}$&\CtwoPoint{Point estimate}&$\CtwoPoint{\bar x}$=$\CtwoPoint{240}$&$\CtwoPoint{240}$ thousand COP; estimated population mean October electricity bill.\\
$\CtwoMoE{\mathrm{MoE}}$&\CtwoMoE{Margin of error}&$\CtwoCritical{z_{1-\alpha/2}}\frac{\CtwoSD{\sigma}}{\sqrt{\CtwoSize{n}}}=\CtwoCritical{1.960}\frac{\CtwoSD{60}}{\sqrt{\CtwoSize{100}}}$&$\CtwoMoE{11.76}$ thousand COP; distance from estimate to either endpoint.\\
$\CtwoEndpoint{L}$&\CtwoEndpoint{Lower endpoint}&$\CtwoPoint{\bar x}-\CtwoMoE{\mathrm{MoE}}=\CtwoPoint{240}-\CtwoMoE{11.76}$&$\CtwoEndpoint{228.24}$ thousand COP; lower bound for the population mean.\\
$\CtwoEndpoint{U}$&\CtwoEndpoint{Upper endpoint}&$\CtwoPoint{\bar x}+\CtwoMoE{\mathrm{MoE}}=\CtwoPoint{240}+\CtwoMoE{11.76}$&$\CtwoEndpoint{251.76}$ thousand COP; upper bound for the population mean.\\
$\CtwoCI{\mathrm{CI}}$&\CtwoCI{Confidence interval}&$[\CtwoEndpoint{L},\CtwoEndpoint{U}]=\CtwoCI{[228.24,251.76]}$&Plausible population mean October electricity bill values, in thousand COP.\\
$\CtwoWidth{W}$&\CtwoWidth{Full width}&$\CtwoEndpoint{U}-\CtwoEndpoint{L}=\CtwoEndpoint{251.76}-\CtwoEndpoint{228.24}$
$=2(\CtwoMoE{11.76})=\CtwoWidth{23.52}$&$\CtwoWidth{23.52}$ thousand COP; smaller width means greater precision.\\}



We are \CtwoConfidence{95\%} confident that the mean October electricity bill for all 200,000 registered retail businesses operating in Medellin throughout October 2026 lies between $\CtwoEndpoint{228.24}$ and $\CtwoEndpoint{251.76}$ thousand COP, in the repeated-sampling sense described in Task 3. This estimates a population mean, not a range containing \CtwoConfidence{95\%} of individual values.

\secInicio[][sec:p1]

\clearpage

\taskprompt{C2 Task 3}{1}{\cTwoC}

\worked{C2 Task 3 \score{1}}{sol:p1c2c}

\textbf{Relevant problem information}

\begin{probleminfoexcerpt}
An energy-cost review concerns all \Pop{200,000} \Pop{registered} \Pop{retail} \Pop{businesses} operating in Medellin throughout October 2026. Each business contributes its total October electricity bill, measured in thousand COP.
\end{probleminfoexcerpt}

\begin{probleminfoexcerpt}
A certified spread-only population report makes the \CtwoSD{population} \CtwoSD{standard} \CtwoSD{deviation} of October bills known as \CtwoSD{60} \CtwoSD{thousand} \CtwoSD{COP}; it does not disclose the \CtwoMean{population} \CtwoMean{mean}.
\end{probleminfoexcerpt}

\begin{probleminfoexcerpt}
Assume the bill distribution is sufficiently regular, without dominant outliers, for a normal approximation to sample means of this size. Each study draws a simple random sample without replacement from the full fixed register. Samples are selected independently across studies; overlap between different samples is permitted. Each study selects only 0.05\% of the population, so dependence within a draw and the finite-population correction are negligible.
\end{probleminfoexcerpt}

\begin{probleminfoexcerpt}
Ten subsequent independent studies each use \CtwoSize{100} businesses, the same known \CtwoSD{population} \CtwoSD{standard} \CtwoSD{deviation} and the same \CtwoConfidence{95\%} normal \CtwoConfidence{confidence} procedure. These are additional studies, separate from both samples above.
\end{probleminfoexcerpt}

\begin{probleminfoexcerpt}
A later comprehensive investigation of the entire October population establishes its \CtwoMean{mean} \CtwoMean{bill} as \CtwoMean{245} \CtwoMean{thousand} \CtwoMean{COP}. Repeated samples do not establish this value. The following subsequent intervals allow retrospective assessment; all endpoints are in thousand COP.
\end{probleminfoexcerpt}

\par\noindent\begin{minipage}{\linewidth}
\begin{minipage}[t]{.49\linewidth}
\keytable{@{}>{\hsize=1.3\hsize}K@{\hspace{2pt}}>{\hsize=.85\hsize}K@{\hspace{2pt}}>{\hsize=.85\hsize}K@{}}{
\textbf{Subsequent study}&\textbf{\CtwoEndpoint{Lower} \CtwoEndpoint{endpoint}}&\textbf{\CtwoEndpoint{Upper} \CtwoEndpoint{endpoint}}\\\midrule
1&\CtwoEndpoint{224.24}&\CtwoEndpoint{247.76}\\
2&\CtwoEndpoint{232.24}&\CtwoEndpoint{255.76}\\
3&\CtwoEndpoint{238.24}&\CtwoEndpoint{261.76}\\
4&\CtwoEndpoint{246.24}&\CtwoEndpoint{269.76}\\
5&\CtwoEndpoint{218.24}&\CtwoEndpoint{241.76}\\
}
\end{minipage}\hfill
\begin{minipage}[t]{.49\linewidth}
\keytable{@{}>{\hsize=1.3\hsize}K@{\hspace{2pt}}>{\hsize=.85\hsize}K@{\hspace{2pt}}>{\hsize=.85\hsize}K@{}}{
\textbf{Subsequent study}&\textbf{\CtwoEndpoint{Lower} \CtwoEndpoint{endpoint}}&\textbf{\CtwoEndpoint{Upper} \CtwoEndpoint{endpoint}}\\\midrule
6&\CtwoEndpoint{235.24}&\CtwoEndpoint{258.76}\\
7&\CtwoEndpoint{230.24}&\CtwoEndpoint{253.76}\\
8&\CtwoEndpoint{242.24}&\CtwoEndpoint{265.76}\\
9&\CtwoEndpoint{221.24}&\CtwoEndpoint{244.76}\\
10&\CtwoEndpoint{234.24}&\CtwoEndpoint{257.76}\\
}
\end{minipage}
\end{minipage}\par

Test each subsequent interval against $\CtwoMean{\mu_{\mathrm{known}}}=\CtwoMean{245}$ thousand COP; include the mean when it equals an endpoint. All endpoints below use thousand COP.

\keytable{KKKK}{
\textbf{Study}&\textbf{Lower endpoint}&\textbf{Upper endpoint}&\textbf{Covers later mean?}\\\midrule
1&$\CtwoEndpoint{224.24}$&$\CtwoEndpoint{247.76}$&\CtwoCoverage{Yes}\\
2&$\CtwoEndpoint{232.24}$&$\CtwoEndpoint{255.76}$&\CtwoCoverage{Yes}\\
3&$\CtwoEndpoint{238.24}$&$\CtwoEndpoint{261.76}$&\CtwoCoverage{Yes}\\
4&$\CtwoEndpoint{246.24}$&$\CtwoEndpoint{269.76}$&\CtwoCoverage{No}\\
5&$\CtwoEndpoint{218.24}$&$\CtwoEndpoint{241.76}$&\CtwoCoverage{No}\\
6&$\CtwoEndpoint{235.24}$&$\CtwoEndpoint{258.76}$&\CtwoCoverage{Yes}\\
7&$\CtwoEndpoint{230.24}$&$\CtwoEndpoint{253.76}$&\CtwoCoverage{Yes}\\
8&$\CtwoEndpoint{242.24}$&$\CtwoEndpoint{265.76}$&\CtwoCoverage{Yes}\\
9&$\CtwoEndpoint{221.24}$&$\CtwoEndpoint{244.76}$&\CtwoCoverage{No}\\
10&$\CtwoEndpoint{234.24}$&$\CtwoEndpoint{257.76}$&\CtwoCoverage{Yes}\\
}


At \CtwoConfidence{95\%}: $\CtwoCoverage{N_{\mathrm{cover}}}=7$ and $N(95\%)=10$ subsequent intervals. Thus
\[\begin{aligned}
\CtwoCoverage{\mathrm{Cov}}(\CtwoConfidence{95\%})&=\frac{\CtwoCoverage{7}}{\CtwoCoverage{10}}\,100\%=\CtwoCoverage{70\%},\\
\CtwoError{E_{\mathrm{abs}}}&=|\CtwoCoverage{70}-\CtwoConfidence{95}|=\CtwoError{25}\text{ percentage points}.
\end{aligned}\]

The \CtwoCoverage{observed coverage} of the later-established mean October electricity bill is \CtwoError{25} percentage points below nominal \CtwoConfidence{95\%}. Each listed interval has full width $2(\CtwoMoE{11.76})=\CtwoWidth{23.52}$ thousand COP, consistent with the stated $n$, known $\CtwoSD{\sigma}$ and critical value. Its center is its own subsequent sample mean, found as $(L_i+U_i)/2$, not the initial mean.

\Needspace{14\baselineskip}
\textbf{Interpretation.} Under the stated model, repeated independent use of each confidence procedure covers the fixed population mean in approximately its nominal percentage of intervals over the long run. A finite collection of intervals need not have exactly nominal coverage, and successive observed proportions need not approach it monotonically. The discrepancies here can arise from finite-study variation; they do not by themselves invalidate the procedure or prove that nominal coverage is wrong.

The population mean is fixed. Once an interval is calculated, it either contains that mean or does not; no nominal probability is assigned to this already determined event. Coverage counts compare each subsequent study's own interval with the later-established population mean. They do not count sample means or individual observations inside the initial interval. The initial and comparison samples are excluded from the subsequent-study denominator.

\secInicio[][sec:p1]

\clearpage

\subsection*{Proposed partial-credit guidance}
\keytable{Kp{.12\linewidth}K}{
\textbf{Task}&\textbf{Points}&\textbf{Evidence earning credit}\\\midrule
C1 Task 1&2&0.5 population and both sample scopes/sizes; 0.75 mean/SD pairings, units and known status; 0.75 counts, both proportions and meanings, unknown population proportion.\\
C1 Task 2&2&1.5 for three signed differences (0.5 each, including percentage points); 0.5 fixed parameters and scope with sampling-variation explanation.\\
C2 Task 1&1&0.25 parameter group and initially unknown/later mean; 0.25 initial sample entries and correct SE with units; 0.5 complete procedure settings at all requested levels.\\
C2 Task 2&2&0.4 point estimate; 0.4 correct margins using known SD; 0.6 endpoints and intervals; 0.3 full widths; 0.3 units and contextual mean interpretation.\\
C2 Task 3&1&0.25 correct covering counts; 0.25 coverage percentages; 0.25 absolute differences in percentage points; 0.25 valid coverage/finite-study interpretation.\\}
\textbf{Task totals:} C1: $2+2=4$; C2: $1+2+1=4$; problem total: 8. Award proportional credit within each listed component for incomplete but valid reasoning. Accept equivalent notation; retain all endpoint and width digits generated by the supplied critical values. Carry forward an earlier numerical error when the later method is correct; do not penalize the same arithmetic error repeatedly. Do not award the interpretation component for a probability statement about the fixed mean in an already calculated interval.


\secInicio[][sec:p1]

\clearpage

\sectionScored{Savings-account interest}{C1-\sml{4}\quad C2-\sml{4}}{sec:p2}

\textbf{Original shared description}\score{8 overall}

A bank reviews all 500,000 positive-balance personal savings accounts at its Cartagena branch network on October 31, 2026. The numerical measurement is the total interest credited to each eligible account during October 2026, in thousand COP. A separate yes-or-no condition records whether the account had mobile banking activated on October 31.

The initial independent random sample contains 144 eligible accounts, with mean October interest of 82 thousand COP, sample standard deviation of 33 thousand COP and 54 accounts with mobile banking activated. A second independent random sample of 144 accounts from exactly the same eligible population and period has mean interest of 79 thousand COP, sample standard deviation of 39 thousand COP and 72 accounts with mobile banking activated.

A certified spread-only report establishes the population standard deviation of October interest as 36 thousand COP while withholding the population mean. This known spread differs from both sample standard deviations. The initial study estimates the initially unknown population mean at 90\% and 99\% confidence.

Assume a sufficiently regular interest distribution without dominant outliers and a valid normal approximation for sample means of this size. Each study draws a simple random sample without replacement from the same full fixed register; draws across studies are independent and may overlap. Each sample contains only 0.0288\% of the population, so the finite-population correction is negligible.

Two hundred subsequent independent studies, separate from the two samples above, each select 144 accounts and use the same known population standard deviation and both confidence procedures. Each study builds both intervals around its own sample mean. The 99\% interval is wider and contains that study's 90\% interval. A later comprehensive investigation of all eligible accounts establishes mean October interest as 84 thousand COP; the repeated samples do not establish it. The table records mutually exclusive retrospective coverage categories for these same 200 studies.

\keytable{Kr}{
\textbf{Intervals covering the later established mean}&\textbf{Studies}\\\midrule
Both confidence levels&176\\
Only the 99\% level&20\\
Neither confidence level&4\\}

Complete all five \textbf{General tasks for every problem} above; see the shared table for task marks.

\textbf{Linked tasks.} Each full task is reproduced immediately before its solution.

\begin{enumerate}[label=(\alph*),itemsep=6pt]
\item \hyperref[sol:p2c1a]{\textbf{C1 Task 1}}. \cOneA\ \sscore{\sml{2}}
\item \hyperref[sol:p2c1b]{\textbf{C1 Task 2}}. \cOneB\ \sscore{\sml{2}}
\item \hyperref[sol:p2c2a]{\textbf{C2 Task 1}}. \cTwoA\ \sscore{\sml{1}}
\item \hyperref[sol:p2c2b]{\textbf{C2 Task 2}}. \cTwoB\ \sscore{\sml{2}}
\item \hyperref[sol:p2c2c]{\textbf{C2 Task 3}}. \cTwoC\ \sscore{\sml{1}}
\end{enumerate}
\secInicio[][sec:p2]

\clearpage

\taskprompt{C1 Task 1}{2}{\cOneA}

\worked{C1 Task 1 \score{2}}{sol:p2c1a}

\textbf{Relevant problem information}

\begin{probleminfoexcerpt}
A bank reviews all \Pop{500,000} \Pop{positive-balance} \Pop{personal} \Pop{savings} \Pop{accounts} at its Cartagena branch network on October 31, 2026. The numerical measurement is the total interest credited to each eligible account during October 2026, in thousand COP. A separate yes-or-no condition records whether the account had mobile banking activated on October 31.
\end{probleminfoexcerpt}

\begin{probleminfoexcerpt}
The \Sam{initial} \Sam{independent} \Sam{random} \Sam{sample} contains \Sam{144} eligible accounts, with \CtwoPoint{mean} \CtwoPoint{October} \CtwoPoint{interest} of \CtwoPoint{82} \CtwoPoint{thousand} \CtwoPoint{COP}, \SSD{sample} \SSD{standard} \SSD{deviation} of \SSD{33} \SSD{thousand} \SSD{COP} and \Sam{54} accounts with mobile banking activated. A \Sam{second} \Sam{independent} \Sam{random} \Sam{sample} of \Sam{144} accounts from exactly the same eligible population and period has \CtwoPoint{mean} \CtwoPoint{interest} of \CtwoPoint{79} \CtwoPoint{thousand} \CtwoPoint{COP}, \SSD{sample} \SSD{standard} \SSD{deviation} of \SSD{39} \SSD{thousand} \SSD{COP} and \Sam{72} accounts with mobile banking activated.
\end{probleminfoexcerpt}

\begin{probleminfoexcerpt}
A certified spread-only report establishes the \CtwoSD{population} \CtwoSD{standard} \CtwoSD{deviation} of \CtwoSD{October interest} as \CtwoSD{36} \CtwoSD{thousand} \CtwoSD{COP} while withholding the \CtwoMean{population} \CtwoMean{mean October} \CtwoMean{interest credit}.
\end{probleminfoexcerpt}

\textbf{Later investigation (retrospective information).}

\begin{probleminfoexcerpt}
A later comprehensive investigation of all eligible accounts establishes \CtwoMean{mean October} \CtwoMean{interest credit} as \CtwoMean{84} \CtwoMean{thousand} \CtwoMean{COP}; the repeated samples do not establish it.
\end{probleminfoexcerpt}

They are separate independent selections, not populations at two different dates.

\textbf{Paired population and sample measures.} All monetary measures below use thousand COP. Sample proportions describe the condition defined in the population table.
\keytable{p{.12\linewidth}p{.24\linewidth}K}{
\textbf{Symbol}&\textbf{Concept}&\textbf{Meaning in context}\\\midrule
$\Pop{\mathcal P,\ N}$&Population and size&\Pop{500,000 eligible accounts} in the scope stated above.\\
$\CtwoMean{\mu}$&\CtwoMean{Population mean}&\CtwoMean{mean October} \CtwoMean{interest credit}; initially unknown, later $\CtwoMean{84}$ thousand COP.\\
$\CtwoSD{\sigma}$&\CtwoSD{Population SD}&$\CtwoSD{36}$ thousand COP; known \CtwoSD{population standard} \CtwoSD{deviation}, the \CtwoSD{spread of October} \CtwoSD{interest credits}.\\
$\PP{p}$&\PP{Population proportion}&\PP{proportion of all} \PP{eligible accounts} \PP{that had mobile banking} \PP{activated on October 31}; unknown.\\}
\keytable{p{.16\linewidth}p{.23\linewidth}KK}{
\textbf{Symbol}&\textbf{Concept}&\textbf{Initial sample A}&\textbf{Second sample B}\\\midrule
$\Sam{X_A,X_B;\ n_A,n_B}$&Samples and sizes&$\Sam{144}$ randomly selected accounts&$\Sam{144}$ independently selected accounts\\
$\CtwoPoint{\bar x_A,\bar x_B}$&\CtwoPoint{Sample mean}&$\CtwoPoint{82}$ thousand COP&$\CtwoPoint{79}$ thousand COP\\
$\SSD{s_A,s_B}$&\SSD{Sample SD}&$\SSD{33}$ thousand COP&$\SSD{39}$ thousand COP\\
$\Sam{k_A,k_B}$&Condition count&$\Sam{54}$ accounts&$\Sam{72}$ accounts\\
$\SP{\hat p_A,\hat p_B}$&\SP{Sample proportion}&$\SP{\hat p_A}=\Sam{54}/\Sam{144}=\SP{0.375=37.5\%}$&$\SP{\hat p_B}=\Sam{72}/\Sam{144}=\SP{0.5=50\%}$\\}



The condition counts are $\Sam{k_A}=\Sam{54}$ and $\Sam{k_B}=\Sam{72}$. The proportions are unitless shares of each sample meeting the stated condition. The known \CtwoSD{population SD} does not make either \SSD{sample SD} equal to it. Initially $\CtwoMean{\mu}$ is unknown; the later investigation establishes that same parameter as $\CtwoMean{84}$ thousand COP. The supplied information never establishes $\PP{p}$: sample proportions are estimates, and the investigation reports only the population mean.

\secInicio[][sec:p2]

\clearpage

\taskprompt{C1 Task 2}{2}{\cOneB}

\worked{C1 Task 2 \score{2}}{sol:p2c1b}

\textbf{Relevant problem information}

\begin{probleminfoexcerpt}
A bank reviews all \Pop{500,000} \Pop{positive-balance} \Pop{personal} \Pop{savings} \Pop{accounts} at its Cartagena branch network on October 31, 2026. The numerical measurement is the total interest credited to each eligible account during October 2026, in thousand COP. A separate yes-or-no condition records whether the account had mobile banking activated on October 31.
\end{probleminfoexcerpt}

\begin{probleminfoexcerpt}
The \Sam{initial} \Sam{independent} \Sam{random} \Sam{sample} contains \Sam{144} eligible accounts, with \CtwoPoint{mean} \CtwoPoint{October} \CtwoPoint{interest} of \CtwoPoint{82} \CtwoPoint{thousand} \CtwoPoint{COP}, \SSD{sample} \SSD{standard} \SSD{deviation} of \SSD{33} \SSD{thousand} \SSD{COP} and \Sam{54} accounts with mobile banking activated. A \Sam{second} \Sam{independent} \Sam{random} \Sam{sample} of \Sam{144} accounts from exactly the same eligible population and period has \CtwoPoint{mean} \CtwoPoint{interest} of \CtwoPoint{79} \CtwoPoint{thousand} \CtwoPoint{COP}, \SSD{sample} \SSD{standard} \SSD{deviation} of \SSD{39} \SSD{thousand} \SSD{COP} and \Sam{72} accounts with mobile banking activated.
\end{probleminfoexcerpt}

\keytable{KKK}{
\textbf{Second minus initial}&\textbf{Calculation}&\textbf{Interpretation}\\\midrule
\CtwoPoint{Mean}&$\CtwoPoint{79-82=-3}$ thousand COP&Second sample mean is lower.\\
\SSD{Sample SD}&$\SSD{39-33=+6}$ thousand COP&Second sample has greater numerical spread.\\
\SP{Proportion}&$\SP{50\%-37.5\%=12.5}$ percentage points&Second sample has a larger condition share.\\}


Both sample sizes remain 144; the selected units and descriptive statistics differ. Population membership, size, place, period, measurement and binary-condition definitions remain fixed. Therefore $\CtwoMean{\mu}$, $\CtwoSD{\sigma}$ and $p$ remain fixed, whether known or unknown. Random sampling variation can produce these differences; they do not establish a temporal change in population parameters. Learning the value of $\CtwoMean{\mu}$ later changes our knowledge, not the population mean itself.

\secInicio[][sec:p2]

\clearpage

\taskprompt{C2 Task 1}{1}{\cTwoA}

\worked{C2 Task 1 \score{1}}{sol:p2c2a}

\textbf{Relevant problem information}

\begin{probleminfoexcerpt}
A bank reviews all \Pop{500,000} \Pop{positive-balance} \Pop{personal} \Pop{savings} \Pop{accounts} at its Cartagena branch network on October 31, 2026. The numerical measurement is the total interest credited to each eligible account during October 2026, in thousand COP.
\end{probleminfoexcerpt}

\begin{probleminfoexcerpt}
The \CtwoSize{initial} \CtwoSize{independent} \CtwoSize{random} \CtwoSize{sample} contains \CtwoSize{144} eligible accounts, with \CtwoPoint{mean} \CtwoPoint{October} \CtwoPoint{interest} of \CtwoPoint{82} \CtwoPoint{thousand} \CtwoPoint{COP}, \SSD{sample} \SSD{standard} \SSD{deviation} of \SSD{33} \SSD{thousand} \SSD{COP} and \Sam{54} accounts with mobile banking activated.
\end{probleminfoexcerpt}

\begin{probleminfoexcerpt}
A certified spread-only report establishes the \CtwoSD{population} \CtwoSD{standard} \CtwoSD{deviation} of \CtwoSD{October interest} as \CtwoSD{36} \CtwoSD{thousand} \CtwoSD{COP} while withholding the \CtwoMean{population} \CtwoMean{mean October} \CtwoMean{interest credit}. The initial study estimates the initially unknown \CtwoMean{population} \CtwoMean{mean October} \CtwoMean{interest credit} at \CtwoConfidence{90\%} and \CtwoConfidence{99\%} \CtwoConfidence{confidence}.
\end{probleminfoexcerpt}

\begin{probleminfoexcerpt}
Assume a sufficiently regular interest distribution without dominant outliers and a valid normal approximation for sample means of this size. Each study draws a simple random sample without replacement from the same full fixed register; draws across studies are independent and may overlap. Each sample contains only 0.0288\% of the population, so the finite-population correction is negligible.
\end{probleminfoexcerpt}

\textbf{Later investigation (unavailable for initial estimation).}

\begin{probleminfoexcerpt}
A later comprehensive investigation of all eligible accounts establishes \CtwoMean{mean October} \CtwoMean{interest credit} as \CtwoMean{84} \CtwoMean{thousand} \CtwoMean{COP}; the repeated samples do not establish it.
\end{probleminfoexcerpt}

\begingroup\raggedright The unknown $\PP{p}$ is the \PP{proportion of all} \PP{eligible accounts} \PP{that had mobile banking} \PP{activated on October 31}. Its denominator is all 500,000 eligible accounts defined above.\par\endgroup

\clearpage
\keytable{p{.12\linewidth}p{.24\linewidth}K}{
\textbf{Symbol}&\textbf{Concept}&\textbf{Meaning in context}\\\midrule
\multicolumn{3}{l}{\textbf{Population parameters}}\\
$\CtwoMean{\mu}$&\CtwoMean{Population mean}&Initially unknown \CtwoMean{mean October} \CtwoMean{interest credit}.\\
$\CtwoMean{\mu_{\mathrm{known}}}$&\CtwoMean{Later established} \CtwoMean{mean}&$\CtwoMean{84}$ thousand COP; the same fixed \CtwoMean{mean October} \CtwoMean{interest credit}, established later by the population investigation.\\
$\CtwoSD{\sigma}$&\CtwoSD{Population SD}&$\CtwoSD{36}$ thousand COP; known \CtwoSD{population standard} \CtwoSD{deviation}, the \CtwoSD{spread of October} \CtwoSD{interest credits}.\\\midrule
\multicolumn{3}{l}{\textbf{Sample statistics}}\\
$\CtwoPoint{\bar x}$&\CtwoPoint{Initial sample mean}&$\CtwoPoint{82}$ thousand COP; point estimate.\\
$\CtwoSize{n}$&\CtwoSize{Initial sample size}&$\CtwoSize{144}$ accounts.\\
$\CtwoSE{\mathrm{SE}}$&\CtwoSE{Standard error}&$\CtwoSE{\mathrm{SE}}=\frac{\CtwoSD{\sigma}}{\sqrt{\CtwoSize{n}}}=\frac{\CtwoSD{36}}{\sqrt{\CtwoSize{144}}}=\CtwoSE{3}$ thousand COP; spread of sample means, using known $\CtwoSD{\sigma}$ and the initial sample size $\CtwoSize{n}$.\\\midrule
\multicolumn{3}{l}{\textbf{Procedure settings}}\\
$\CtwoConfidence{C}$&\CtwoConfidence{Confidence}&$\CtwoConfidence{90\%}$: $\CtwoConfidence{0.9}$; $\CtwoConfidence{99\%}$: $\CtwoConfidence{0.99}$.\\
$\CtwoSignificance{\alpha}$&\CtwoSignificance{Significance}&$\CtwoConfidence{90\%}$: $1-\CtwoConfidence{0.9}=\CtwoSignificance{0.1}$; $\CtwoConfidence{99\%}$: $1-\CtwoConfidence{0.99}=\CtwoSignificance{0.01}$.\\
$\CtwoSignificance{\alpha/2}$&\CtwoSignificance{Each tail area}&$\CtwoConfidence{90\%}$: $\CtwoSignificance{0.1}/2=\CtwoSignificance{0.05}$; $\CtwoConfidence{99\%}$: $\CtwoSignificance{0.01}/2=\CtwoSignificance{0.005}$.\\
$\CtwoCritical{z_{1-\alpha/2}}$&\CtwoCritical{Positive} \CtwoCritical{critical value}&$\CtwoConfidence{90\%}$: $\CtwoCritical{1.645}$; $\CtwoConfidence{99\%}$: $\CtwoCritical{2.576}$.\\}



Use only the initial sample for estimation; the comparison sample describes variation, and the later mean supports retrospective assessment. Use known $\CtwoSD{\sigma}=\CtwoSD{36}$, rather than $\SSD{s_A}=\SSD{33}$ or $\SSD{s_B}=\SSD{39}$. Random sampling, the stated regularity assumption and negligible sample fraction support the normal approximation.

\secInicio[][sec:p2]

\clearpage

\clearpage

\taskprompt{C2 Task 2}{2}{\cTwoB}

\worked{C2 Task 2 \score{2}}{sol:p2c2b}

\textbf{Relevant problem information}

\begin{probleminfoexcerpt}
A bank reviews all \Pop{500,000} \Pop{positive-balance} \Pop{personal} \Pop{savings} \Pop{accounts} at its Cartagena branch network on October 31, 2026. The numerical measurement is the total interest credited to each eligible account during October 2026, in thousand COP.
\end{probleminfoexcerpt}

\begin{probleminfoexcerpt}
The \CtwoSize{initial} \CtwoSize{independent} \CtwoSize{random} \CtwoSize{sample} contains \CtwoSize{144} eligible accounts, with \CtwoPoint{mean} \CtwoPoint{October} \CtwoPoint{interest} of \CtwoPoint{82} \CtwoPoint{thousand} \CtwoPoint{COP}, \SSD{sample} \SSD{standard} \SSD{deviation} of \SSD{33} \SSD{thousand} \SSD{COP} and \Sam{54} accounts with mobile banking activated.
\end{probleminfoexcerpt}

\begin{probleminfoexcerpt}
A certified spread-only report establishes the \CtwoSD{population} \CtwoSD{standard} \CtwoSD{deviation} of October interest as \CtwoSD{36} \CtwoSD{thousand} \CtwoSD{COP} while withholding the \CtwoMean{population} \CtwoMean{mean}. The initial study estimates the initially unknown \CtwoMean{population} \CtwoMean{mean} at \CtwoConfidence{90\%} and \CtwoConfidence{99\%} \CtwoConfidence{confidence}.
\end{probleminfoexcerpt}

\begin{probleminfoexcerpt}
Assume a sufficiently regular interest distribution without dominant outliers and a valid normal approximation for sample means of this size. Each study draws a simple random sample without replacement from the same full fixed register; draws across studies are independent and may overlap. Each sample contains only 0.0288\% of the population, so the finite-population correction is negligible.
\end{probleminfoexcerpt}

\textbf{Later investigation (unavailable for initial estimation).}

\begin{probleminfoexcerpt}
A later comprehensive investigation of all eligible accounts establishes \CtwoMean{mean} \CtwoMean{October} \CtwoMean{interest} as \CtwoMean{84} \CtwoMean{thousand} \CtwoMean{COP}; the repeated samples do not establish it.
\end{probleminfoexcerpt}

The shared \CtwoPoint{point estimate} is $\CtwoPoint{\bar x}=\CtwoPoint{82}$ thousand COP, the average October interest credit in the initial sample; it estimates the population mean with a single value.

\textbf{\CtwoConfidence{90\%} confidence}

\keytable{p{.10\linewidth}p{.17\linewidth}p{.32\linewidth}K}{
\textbf{Symbol}&\textbf{Concept}&\textbf{Calculation}&\textbf{Result and meaning in context}\\\midrule
$\CtwoPoint{\bar x}$&\CtwoPoint{Point estimate}&$\CtwoPoint{\bar x}$=$\CtwoPoint{82}$&$\CtwoPoint{82}$ thousand COP; estimated population mean October interest credit.\\
$\CtwoMoE{\mathrm{MoE}}$&\CtwoMoE{Margin of error}&$\CtwoCritical{z_{1-\alpha/2}}\frac{\CtwoSD{\sigma}}{\sqrt{\CtwoSize{n}}}=\CtwoCritical{1.645}\frac{\CtwoSD{36}}{\sqrt{\CtwoSize{144}}}$&$\CtwoMoE{4.935}$ thousand COP; distance from estimate to either endpoint.\\
$\CtwoEndpoint{L}$&\CtwoEndpoint{Lower endpoint}&$\CtwoPoint{\bar x}-\CtwoMoE{\mathrm{MoE}}=\CtwoPoint{82}-\CtwoMoE{4.935}$&$\CtwoEndpoint{77.065}$ thousand COP; lower bound for the population mean.\\
$\CtwoEndpoint{U}$&\CtwoEndpoint{Upper endpoint}&$\CtwoPoint{\bar x}+\CtwoMoE{\mathrm{MoE}}=\CtwoPoint{82}+\CtwoMoE{4.935}$&$\CtwoEndpoint{86.935}$ thousand COP; upper bound for the population mean.\\
$\CtwoCI{\mathrm{CI}}$&\CtwoCI{Confidence interval}&$[\CtwoEndpoint{L},\CtwoEndpoint{U}]=\CtwoCI{[77.065,86.935]}$&Plausible population mean October interest credit values, in thousand COP.\\
$\CtwoWidth{W}$&\CtwoWidth{Full width}&$\CtwoEndpoint{U}-\CtwoEndpoint{L}=\CtwoEndpoint{86.935}-\CtwoEndpoint{77.065}$
$=2(\CtwoMoE{4.935})=\CtwoWidth{9.87}$&$\CtwoWidth{9.87}$ thousand COP; smaller width means greater precision.\\}



We are \CtwoConfidence{90\%} confident that the mean October interest credit for all 500,000 positive-balance personal savings accounts at the Cartagena branch network on October 31, 2026 lies between $\CtwoEndpoint{77.065}$ and $\CtwoEndpoint{86.935}$ thousand COP, in the repeated-sampling sense described in Task 3. This estimates a population mean, not a range containing \CtwoConfidence{90\%} of individual values.

\Needspace{22\baselineskip}
\textbf{\CtwoConfidence{99\%} confidence}

\keytable{p{.10\linewidth}p{.17\linewidth}p{.32\linewidth}K}{
\textbf{Symbol}&\textbf{Concept}&\textbf{Calculation}&\textbf{Result and meaning in context}\\\midrule
$\CtwoPoint{\bar x}$&\CtwoPoint{Point estimate}&$\CtwoPoint{\bar x}$=$\CtwoPoint{82}$&$\CtwoPoint{82}$ thousand COP; estimated population mean October interest credit.\\
$\CtwoMoE{\mathrm{MoE}}$&\CtwoMoE{Margin of error}&$\CtwoCritical{z_{1-\alpha/2}}\frac{\CtwoSD{\sigma}}{\sqrt{\CtwoSize{n}}}=\CtwoCritical{2.576}\frac{\CtwoSD{36}}{\sqrt{\CtwoSize{144}}}$&$\CtwoMoE{7.728}$ thousand COP; distance from estimate to either endpoint.\\
$\CtwoEndpoint{L}$&\CtwoEndpoint{Lower endpoint}&$\CtwoPoint{\bar x}-\CtwoMoE{\mathrm{MoE}}=\CtwoPoint{82}-\CtwoMoE{7.728}$&$\CtwoEndpoint{74.272}$ thousand COP; lower bound for the population mean.\\
$\CtwoEndpoint{U}$&\CtwoEndpoint{Upper endpoint}&$\CtwoPoint{\bar x}+\CtwoMoE{\mathrm{MoE}}=\CtwoPoint{82}+\CtwoMoE{7.728}$&$\CtwoEndpoint{89.728}$ thousand COP; upper bound for the population mean.\\
$\CtwoCI{\mathrm{CI}}$&\CtwoCI{Confidence interval}&$[\CtwoEndpoint{L},\CtwoEndpoint{U}]=\CtwoCI{[74.272,89.728]}$&Plausible population mean October interest credit values, in thousand COP.\\
$\CtwoWidth{W}$&\CtwoWidth{Full width}&$\CtwoEndpoint{U}-\CtwoEndpoint{L}=\CtwoEndpoint{89.728}-\CtwoEndpoint{74.272}$
$=2(\CtwoMoE{7.728})=\CtwoWidth{15.456}$&$\CtwoWidth{15.456}$ thousand COP; smaller width means greater precision.\\}



We are \CtwoConfidence{99\%} confident that the mean October interest credit for all 500,000 positive-balance personal savings accounts at the Cartagena branch network on October 31, 2026 lies between $\CtwoEndpoint{74.272}$ and $\CtwoEndpoint{89.728}$ thousand COP, in the repeated-sampling sense described in Task 3. This estimates a population mean, not a range containing \CtwoConfidence{99\%} of individual values.

The \CtwoConfidence{99\%} interval contains the \CtwoConfidence{90\%} interval: the same center and SE are used, and the larger critical value produces a wider interval and less precision at the higher confidence level.

\secInicio[][sec:p2]

\clearpage

\taskprompt{C2 Task 3}{1}{\cTwoC}

\worked{C2 Task 3 \score{1}}{sol:p2c2c}

\textbf{Relevant problem information}

\begin{probleminfoexcerpt}
A bank reviews all \Pop{500,000} \Pop{positive-balance} \Pop{personal} \Pop{savings} \Pop{accounts} at its Cartagena branch network on October 31, 2026. The numerical measurement is the total interest credited to each eligible account during October 2026, in thousand COP.
\end{probleminfoexcerpt}

\begin{probleminfoexcerpt}
A certified spread-only report establishes the \CtwoSD{population} \CtwoSD{standard} \CtwoSD{deviation} of October interest as \CtwoSD{36} \CtwoSD{thousand} \CtwoSD{COP} while withholding the \CtwoMean{population} \CtwoMean{mean}. The initial study estimates the initially unknown \CtwoMean{population} \CtwoMean{mean} at \CtwoConfidence{90\%} and \CtwoConfidence{99\%} \CtwoConfidence{confidence}.
\end{probleminfoexcerpt}

\begin{probleminfoexcerpt}
Assume a sufficiently regular interest distribution without dominant outliers and a valid normal approximation for sample means of this size. Each study draws a simple random sample without replacement from the same full fixed register; draws across studies are independent and may overlap. Each sample contains only 0.0288\% of the population, so the finite-population correction is negligible.
\end{probleminfoexcerpt}

\begin{probleminfoexcerpt}
Two hundred subsequent independent studies, separate from the two samples above, each select \CtwoSize{144} accounts and use the same known \CtwoSD{population} \CtwoSD{standard} \CtwoSD{deviation} and both \CtwoConfidence{confidence} procedures. Each study builds both intervals around its own sample mean. The \CtwoConfidence{99\%} interval is wider and contains that study's \CtwoConfidence{90\%} interval.
\end{probleminfoexcerpt}

\begin{probleminfoexcerpt}
A later comprehensive investigation of all eligible accounts establishes \CtwoMean{mean} \CtwoMean{October} \CtwoMean{interest} as \CtwoMean{84} \CtwoMean{thousand} \CtwoMean{COP}; the repeated samples do not establish it. The table records mutually exclusive retrospective coverage categories for these same 200 studies.
\end{probleminfoexcerpt}

\keytable{Kr}{
\textbf{Intervals covering the later established mean}&\textbf{Studies}\\\midrule
Both confidence levels&\CtwoCoverage{176}\\
Only the \CtwoConfidence{99\%} level&\CtwoCoverage{20}\\
Neither confidence level&\CtwoCoverage{4}\\}

\keytable{KKK}{
\textbf{Confidence}&\textbf{Covering and missing counts}&\textbf{Coverage and absolute difference}\\\midrule
$\CtwoConfidence{90\%}$&Cover: $\CtwoCoverage{176}$; miss: $\CtwoCoverage{20+4=24}$&$\CtwoCoverage{176}/\CtwoCoverage{200}\times100\%=\CtwoCoverage{88\%}$; $|\CtwoCoverage{88}-\CtwoConfidence{90}|=\CtwoError{2}$ percentage points.\\
$\CtwoConfidence{99\%}$&Cover: $\CtwoCoverage{176+20=196}$; miss: $\CtwoCoverage{4}$&$\CtwoCoverage{196}/\CtwoCoverage{200}\times100\%=\CtwoCoverage{98\%}$; $|\CtwoCoverage{98}-\CtwoConfidence{99}|=\CtwoError{1}$ percentage point.\\}


The 90\% procedure covers the fixed mean of 84 thousand COP in 176 of 200 subsequent studies, 2 percentage points below \CtwoConfidence{nominal confidence}. The 99\% procedure covers it in 196 of those same studies, 1 percentage point below \CtwoConfidence{nominal confidence}. The categories sum to $176+20+4=200$. Every 90\% covering interval also covers at 99\%; all four misses at 99\% are among the 24 misses at 90\%. There is no impossible category in which only the narrower interval covers. These are feasible counts: sample-mean distances from 84 of 0, 6 and 9 thousand COP respectively would realize the three categories, since the margins are 4.935 and 7.728.

\Needspace{14\baselineskip}
\textbf{Interpretation.} Under the stated model, repeated independent use of each confidence procedure covers the fixed population mean in approximately its nominal percentage of intervals over the long run. A finite collection of intervals need not have exactly nominal coverage, and successive observed proportions need not approach it monotonically. The discrepancies here can arise from finite-study variation; they do not by themselves invalidate the procedure or prove that nominal coverage is wrong.

The population mean is fixed. Once an interval is calculated, it either contains that mean or does not; no nominal probability is assigned to this already determined event. Coverage counts compare each subsequent study's own interval with the later-established population mean. They do not count sample means or individual observations inside the initial interval. The initial and comparison samples are excluded from the subsequent-study denominator.

\secInicio[][sec:p2]

\clearpage

\subsection*{Proposed partial-credit guidance}
\keytable{Kp{.12\linewidth}K}{
\textbf{Task}&\textbf{Points}&\textbf{Evidence earning credit}\\\midrule
C1 Task 1&2&0.5 population and both sample scopes/sizes; 0.75 mean/SD pairings, units and known status; 0.75 counts, both proportions and meanings, unknown population proportion.\\
C1 Task 2&2&1.5 for three signed differences (0.5 each, including percentage points); 0.5 fixed parameters and scope with sampling-variation explanation.\\
C2 Task 1&1&0.25 parameter group and initially unknown/later mean; 0.25 initial sample entries and correct SE with units; 0.5 complete procedure settings at all requested levels.\\
C2 Task 2&2&0.4 point estimate; 0.4 correct margins using known SD; 0.6 endpoints and intervals; 0.3 full widths; 0.3 units and contextual mean interpretation. Divide each multi-level component equally across the two levels in Problem 2.\\
C2 Task 3&1&0.25 correct covering counts; 0.25 coverage percentages; 0.25 absolute differences in percentage points; 0.25 valid coverage/finite-study interpretation. Divide numerical components equally across levels in Problem 2.\\}
\textbf{Task totals:} C1: $2+2=4$; C2: $1+2+1=4$; problem total: 8. Award proportional credit within each listed component for incomplete but valid reasoning. Accept equivalent notation; retain all endpoint and width digits generated by the supplied critical values. Carry forward an earlier numerical error when the later method is correct; do not penalize the same arithmetic error repeatedly. Do not award the interpretation component for a probability statement about the fixed mean in an already calculated interval.


\secInicio[][sec:p2]

\clearpage

\sectionScored{Create an integrated scenario from scratch}{C1-\sml{6}\quad C2-\sml{6}}{sec:p3}

\textbf{Original construction question}\score{12 overall}
Invent a new economic or financial scenario integrating both criteria. First supply a complete, valid economic or financial setting, including the shared contextual description and all data required by the checklist. Then complete the five tasks below for your own scenario. One confidence level is sufficient. Write the contextual description in ordinary language: numerals, percentages, written units and plain-language table headings are allowed, but reserve symbols, formulas, inequalities and interval notation for task reference material and solutions.

\textbf{Required construction checklist}
\begin{enumerate}
\item Define a population's place, fixed period and size; define one numerical measurement with units and a separate binary condition.
\item Provide an initial independent random sample and a second independent random sample from exactly the same population and period. For each give size, mean, sample standard deviation and integer condition count between zero and that sample size. Ensure summaries are feasible for the numerical measurement.
\item State a known population standard deviation, distinguish it from both sample standard deviations, and explain that the population mean is initially unknown. Choose confidence and a supplied standard normal critical value.
\item Give suitable randomness, independence and normal-approximation assumptions. For finite sampling without replacement, keep the sampling fractions negligible. Subsequent studies must be independent draws from the same full fixed population and period.
\item Specify a number of subsequent studies, each with the initial sample size, known population standard deviation and confidence procedure. Supply all subsequent endpoints or complete coverage/miss counts. Endpoints must have the stated procedure's width; with multiple levels, evidence must respect nesting.
\item Include a later comprehensive population investigation establishing the mean of that same population for retrospective coverage. A clearly hypothetical investigation is acceptable; do not claim repeated samples establish the population mean.
\end{enumerate}

Scenario completeness and validity are assessed within the relevant task marks in the shared table above; no separate construction marks are added.

Complete all five \textbf{General tasks for every problem} above; see the shared table for task marks.

\textbf{Tasks to complete for your scenario.} Follow the five linked tasks below. Other internally consistent scenarios and correct solutions are equally acceptable; the example after the task list is one possible response.

\clearpage
\textbf{Linked tasks.} Each full task is reproduced immediately before its solution.

\begin{enumerate}[label=(\alph*),itemsep=6pt]
\item \hyperref[sol:p3c1a]{\textbf{C1 Task 1}}. \cOneA\ \sscore{\sml{3}}
\item \hyperref[sol:p3c1b]{\textbf{C1 Task 2}}. \cOneB\ \sscore{\sml{3}}
\item \hyperref[sol:p3c2a]{\textbf{C2 Task 1}}. \cTwoA\ \sscore{\sml{2}}
\item \hyperref[sol:p3c2b]{\textbf{C2 Task 2}}. \cTwoB\ \sscore{\sml{2}}
\item \hyperref[sol:p3c2c]{\textbf{C2 Task 3}}. \cTwoC\ \sscore{\sml{2}}
\end{enumerate}
\secInicio[][sec:p3]

\clearpage

\textbf{One valid example of a student-created scenario (hypothetical).}
A platform investigates all 100,000 completed deliveries by its registered couriers in Bogota on November 2, 2026. The measurement is the fee for each delivery in thousand COP. The separate yes-or-no condition is whether the delivery included optional parcel insurance.

The initial independent random sample of 64 deliveries has mean fee 32 thousand COP, sample standard deviation 7 thousand COP and 24 insured deliveries. A second independent random sample of 64 deliveries from that identical population and date has mean fee 34 thousand COP, sample standard deviation 9 thousand COP and 32 insured deliveries.

A certified spread-only report gives the known population standard deviation of delivery fees as 8 thousand COP, distinct from either sample standard deviation; the population mean is initially withheld. The initial study uses 95\% confidence and the supplied positive standard normal critical value 1.960. Assume the fee distribution is sufficiently regular with no dominant outliers, so the normal approximation to sample means is suitable. Each sample is selected without replacement from the full register. Study selections are independent and may overlap. Each study samples 0.064\% of the population, so the finite-population correction is negligible.

Five subsequent independent studies, separate from the two original samples, use 64 deliveries each, the same known population standard deviation and the same 95\% normal procedure. A hypothetical later comprehensive investigation of the entire same-day delivery population establishes the population mean as 33 thousand COP; this is not established by the repeated samples. The subsequent interval endpoints below are in thousand COP.

\keytable{KKK}{
\textbf{Subsequent study}&\textbf{Lower endpoint}&\textbf{Upper endpoint}\\\midrule
1&30.04&33.96\\
2&31.04&34.96\\
3&32.04&35.96\\
4&33.04&36.96\\
5&30.54&34.46\\
}


\textbf{Validity of the construction.} Counts are integers within sample sizes. The numerical fee is nonnegative; the reported means and spreads are feasible. For example, a 64-value sample with 32 values at $m-d$ and 32 at $m+d$, where $d=s\sqrt{63/64}$, has exactly mean $m$ and sample SD $s$. For the first sample $d\approx6.9451<32$; for the second $d\approx8.9294<34$, so all fees can be positive. Insurance labels can be assigned to 24 and 32 of those observations independently of fee size. The large remainder of the population can accommodate both samples and the stated mean and spread. This construction illustrates feasibility; students need not submit raw observations.

\secInicio[][sec:p3]

\clearpage

\taskprompt{C1 Task 1}{3}{\cOneA}

\worked{C1 Task 1 \score{3}}{sol:p3c1a}

\textbf{Relevant problem information}

\begin{probleminfoexcerpt}
A platform investigates all \Pop{100,000} \Pop{completed} \Pop{deliveries} by its registered couriers in Bogota on November 2, 2026. The measurement is the fee for each delivery in thousand COP. The separate yes-or-no condition is whether the delivery included optional parcel insurance.
\end{probleminfoexcerpt}

\begin{probleminfoexcerpt}
The \Sam{initial} \Sam{independent} \Sam{random} \Sam{sample} of \Sam{64} deliveries has \CtwoPoint{mean} \CtwoPoint{fee} \CtwoPoint{32} \CtwoPoint{thousand} \CtwoPoint{COP}, \SSD{sample} \SSD{standard} \SSD{deviation} \SSD{7} \SSD{thousand} \SSD{COP} and \Sam{24} insured deliveries. A \Sam{second} \Sam{independent} \Sam{random} \Sam{sample} of \Sam{64} deliveries from that identical population and date has \CtwoPoint{mean} \CtwoPoint{fee} \CtwoPoint{34} \CtwoPoint{thousand} \CtwoPoint{COP}, \SSD{sample} \SSD{standard} \SSD{deviation} \SSD{9} \SSD{thousand} \SSD{COP} and \Sam{32} insured deliveries.
\end{probleminfoexcerpt}

\begin{probleminfoexcerpt}
A certified spread-only report gives the known \CtwoSD{population} \CtwoSD{standard} \CtwoSD{deviation} of \CtwoSD{delivery fees} as \CtwoSD{8} \CtwoSD{thousand} \CtwoSD{COP}, distinct from either \SSD{sample} \SSD{standard} \SSD{deviation}; the \CtwoMean{population} \CtwoMean{mean delivery fee} is initially withheld.
\end{probleminfoexcerpt}

\textbf{Later investigation (retrospective information).}

\begin{probleminfoexcerpt}
A hypothetical later comprehensive investigation of the entire same-day delivery population establishes the \CtwoMean{population} \CtwoMean{mean delivery fee} as \CtwoMean{33} \CtwoMean{thousand} \CtwoMean{COP}; this is not established by the repeated samples.
\end{probleminfoexcerpt}

They are separate independent selections, not populations at two different dates.

\textbf{Paired population and sample measures.} All monetary measures below use thousand COP. Sample proportions describe the condition defined in the population table.
\keytable{p{.12\linewidth}p{.24\linewidth}K}{
\textbf{Symbol}&\textbf{Concept}&\textbf{Meaning in context}\\\midrule
$\Pop{\mathcal P,\ N}$&Population and size&\Pop{100,000 eligible deliveries} in the scope stated above.\\
$\CtwoMean{\mu}$&\CtwoMean{Population mean}&\CtwoMean{mean delivery fee}; initially unknown, later $\CtwoMean{33}$ thousand COP.\\
$\CtwoSD{\sigma}$&\CtwoSD{Population SD}&$\CtwoSD{8}$ thousand COP; known \CtwoSD{population standard} \CtwoSD{deviation}, the \CtwoSD{spread of delivery fees}.\\
$\PP{p}$&\PP{Population proportion}&\PP{proportion of all} \PP{eligible deliveries} \PP{that included optional} \PP{parcel insurance}; unknown.\\}
\keytable{p{.16\linewidth}p{.23\linewidth}KK}{
\textbf{Symbol}&\textbf{Concept}&\textbf{Initial sample A}&\textbf{Second sample B}\\\midrule
$\Sam{X_A,X_B;\ n_A,n_B}$&Samples and sizes&$\Sam{64}$ randomly selected deliveries&$\Sam{64}$ independently selected deliveries\\
$\CtwoPoint{\bar x_A,\bar x_B}$&\CtwoPoint{Sample mean}&$\CtwoPoint{32}$ thousand COP&$\CtwoPoint{34}$ thousand COP\\
$\SSD{s_A,s_B}$&\SSD{Sample SD}&$\SSD{7}$ thousand COP&$\SSD{9}$ thousand COP\\
$\Sam{k_A,k_B}$&Condition count&$\Sam{24}$ deliveries&$\Sam{32}$ deliveries\\
$\SP{\hat p_A,\hat p_B}$&\SP{Sample proportion}&$\SP{\hat p_A}=\Sam{24}/\Sam{64}=\SP{0.375=37.5\%}$&$\SP{\hat p_B}=\Sam{32}/\Sam{64}=\SP{0.5=50\%}$\\}



The condition counts are $\Sam{k_A}=\Sam{24}$ and $\Sam{k_B}=\Sam{32}$. The proportions are unitless shares of each sample meeting the stated condition. The known \CtwoSD{population SD} does not make either \SSD{sample SD} equal to it. Initially $\CtwoMean{\mu}$ is unknown; the later investigation establishes that same parameter as $\CtwoMean{33}$ thousand COP. The supplied information never establishes $\PP{p}$: sample proportions are estimates, and the investigation reports only the population mean.

\secInicio[][sec:p3]

\clearpage

\taskprompt{C1 Task 2}{3}{\cOneB}

\worked{C1 Task 2 \score{3}}{sol:p3c1b}

\textbf{Relevant problem information}

\begin{probleminfoexcerpt}
A platform investigates all \Pop{100,000} \Pop{completed} \Pop{deliveries} by its registered couriers in Bogota on November 2, 2026. The measurement is the fee for each delivery in thousand COP. The separate yes-or-no condition is whether the delivery included optional parcel insurance.
\end{probleminfoexcerpt}

\begin{probleminfoexcerpt}
The \Sam{initial} \Sam{independent} \Sam{random} \Sam{sample} of \Sam{64} deliveries has \CtwoPoint{mean} \CtwoPoint{fee} \CtwoPoint{32} \CtwoPoint{thousand} \CtwoPoint{COP}, \SSD{sample} \SSD{standard} \SSD{deviation} \SSD{7} \SSD{thousand} \SSD{COP} and \Sam{24} insured deliveries. A \Sam{second} \Sam{independent} \Sam{random} \Sam{sample} of \Sam{64} deliveries from that identical population and date has \CtwoPoint{mean} \CtwoPoint{fee} \CtwoPoint{34} \CtwoPoint{thousand} \CtwoPoint{COP}, \SSD{sample} \SSD{standard} \SSD{deviation} \SSD{9} \SSD{thousand} \SSD{COP} and \Sam{32} insured deliveries.
\end{probleminfoexcerpt}

\keytable{KKK}{
\textbf{Second minus initial}&\textbf{Calculation}&\textbf{Interpretation}\\\midrule
\CtwoPoint{Mean}&$\CtwoPoint{34-32=+2}$ thousand COP&Second sample mean is higher.\\
\SSD{Sample SD}&$\SSD{9-7=+2}$ thousand COP&Second sample has greater numerical spread.\\
\SP{Proportion}&$\SP{50\%-37.5\%=12.5}$ percentage points&Second sample has a larger condition share.\\}


Both sample sizes remain 64; the selected units and descriptive statistics differ. Population membership, size, place, period, measurement and binary-condition definitions remain fixed. Therefore $\CtwoMean{\mu}$, $\CtwoSD{\sigma}$ and $p$ remain fixed, whether known or unknown. Random sampling variation can produce these differences; they do not establish a temporal change in population parameters. Learning the value of $\CtwoMean{\mu}$ later changes our knowledge, not the population mean itself.

\secInicio[][sec:p3]

\clearpage

\taskprompt{C2 Task 1}{2}{\cTwoA}

\worked{C2 Task 1 \score{2}}{sol:p3c2a}

\textbf{Relevant problem information}

\begin{probleminfoexcerpt}
A platform investigates all \Pop{100,000} \Pop{completed} \Pop{deliveries} by its registered couriers in Bogota on November 2, 2026. The measurement is the fee for each delivery in thousand COP.
\end{probleminfoexcerpt}

\begin{probleminfoexcerpt}
The \CtwoSize{initial} \CtwoSize{independent} \CtwoSize{random} \CtwoSize{sample} of \CtwoSize{64} deliveries has \CtwoPoint{mean} \CtwoPoint{fee} \CtwoPoint{32} \CtwoPoint{thousand} \CtwoPoint{COP}, \SSD{sample} \SSD{standard} \SSD{deviation} \SSD{7} \SSD{thousand} \SSD{COP} and \Sam{24} insured deliveries.
\end{probleminfoexcerpt}

\begin{probleminfoexcerpt}
A certified spread-only report gives the known \CtwoSD{population} \CtwoSD{standard} \CtwoSD{deviation} of \CtwoSD{delivery fees} as \CtwoSD{8} \CtwoSD{thousand} \CtwoSD{COP}, distinct from either \SSD{sample} \SSD{standard} \SSD{deviation}; the \CtwoMean{population} \CtwoMean{mean delivery fee} is initially withheld. The initial study uses \CtwoConfidence{95\%} \CtwoConfidence{confidence} and the supplied \CtwoCritical{positive} \CtwoCritical{standard} \CtwoCritical{normal} \CtwoCritical{critical} \CtwoCritical{value} \CtwoCritical{1.960}.
\end{probleminfoexcerpt}

\begin{probleminfoexcerpt}
Assume the fee distribution is sufficiently regular with no dominant outliers, so the normal approximation to sample means is suitable. Each sample is selected without replacement from the full register. Study selections are independent and may overlap. Each study samples 0.064\% of the population, so the finite-population correction is negligible.
\end{probleminfoexcerpt}

\textbf{Later investigation (unavailable for initial estimation).}

\begin{probleminfoexcerpt}
A hypothetical later comprehensive investigation of the entire same-day delivery population establishes the \CtwoMean{population} \CtwoMean{mean delivery fee} as \CtwoMean{33} \CtwoMean{thousand} \CtwoMean{COP}; this is not established by the repeated samples.
\end{probleminfoexcerpt}

\begingroup\raggedright The unknown $\PP{p}$ is the \PP{proportion of all} \PP{eligible deliveries} \PP{that included optional} \PP{parcel insurance}. Its denominator is all 100,000 eligible deliveries defined above.\par\endgroup

\keytable{p{.12\linewidth}p{.24\linewidth}K}{
\textbf{Symbol}&\textbf{Concept}&\textbf{Meaning in context}\\\midrule
\multicolumn{3}{l}{\textbf{Population parameters}}\\
$\CtwoMean{\mu}$&\CtwoMean{Population mean}&Initially unknown \CtwoMean{mean delivery fee}.\\
$\CtwoMean{\mu_{\mathrm{known}}}$&\CtwoMean{Later established} \CtwoMean{mean}&$\CtwoMean{33}$ thousand COP; the same fixed \CtwoMean{mean delivery fee}, established by the population investigation; unavailable for initial estimation.\\
$\CtwoSD{\sigma}$&\CtwoSD{Population SD}&$\CtwoSD{8}$ thousand COP; known \CtwoSD{population standard} \CtwoSD{deviation}, the \CtwoSD{spread of delivery fees}.\\\midrule
\multicolumn{3}{l}{\textbf{Sample statistics}}\\
$\CtwoPoint{\bar x}$&\CtwoPoint{Initial sample mean}&$\CtwoPoint{32}$ thousand COP; point estimate.\\
$\CtwoSize{n}$&\CtwoSize{Initial sample size}&$\CtwoSize{64}$ deliveries.\\
$\CtwoSE{\mathrm{SE}}$&\CtwoSE{Standard error}&$\CtwoSE{\mathrm{SE}}=\frac{\CtwoSD{\sigma}}{\sqrt{\CtwoSize{n}}}=\frac{\CtwoSD{8}}{\sqrt{\CtwoSize{64}}}=\CtwoSE{1}$ thousand COP; spread of sample means, using known $\CtwoSD{\sigma}$ and the initial sample size $\CtwoSize{n}$.\\\midrule
\multicolumn{3}{l}{\textbf{Procedure settings}}\\
$\CtwoConfidence{C}$&\CtwoConfidence{Confidence}&$\CtwoConfidence{95\%}$: $\CtwoConfidence{0.95}$.\\
$\CtwoSignificance{\alpha}$&\CtwoSignificance{Significance}&$\CtwoConfidence{95\%}$: $1-\CtwoConfidence{0.95}=\CtwoSignificance{0.05}$.\\
$\CtwoSignificance{\alpha/2}$&\CtwoSignificance{Each tail area}&$\CtwoConfidence{95\%}$: $\CtwoSignificance{0.05}/2=\CtwoSignificance{0.025}$.\\
$\CtwoCritical{z_{1-\alpha/2}}$&\CtwoCritical{Positive} \CtwoCritical{critical value}&$\CtwoConfidence{95\%}$: $\CtwoCritical{1.960}$.\\}



Use the initial sample only for estimation. The comparison sample describes sample variation; the later mean is reserved for retrospective assessment. Known $\CtwoSD{\sigma}=\CtwoSD{8}$ is used, rather than either $\SSD{s_A}=\SSD{7}$ or $\SSD{s_B}=\SSD{9}$. The normal approximation is supported by random sampling, the stated regularity assumption and negligible sample fraction.

\secInicio[][sec:p3]

\clearpage

\taskprompt{C2 Task 2}{2}{\cTwoB}

\worked{C2 Task 2 \score{2}}{sol:p3c2b}

\textbf{Relevant problem information}

\begin{probleminfoexcerpt}
A platform investigates all \Pop{100,000} \Pop{completed} \Pop{deliveries} by its registered couriers in Bogota on November 2, 2026. The measurement is the fee for each delivery in thousand COP.
\end{probleminfoexcerpt}

\begin{probleminfoexcerpt}
The \CtwoSize{initial} \CtwoSize{independent} \CtwoSize{random} \CtwoSize{sample} of \CtwoSize{64} deliveries has \CtwoPoint{mean} \CtwoPoint{fee} \CtwoPoint{32} \CtwoPoint{thousand} \CtwoPoint{COP}, \SSD{sample} \SSD{standard} \SSD{deviation} \SSD{7} \SSD{thousand} \SSD{COP} and \Sam{24} insured deliveries.
\end{probleminfoexcerpt}

\begin{probleminfoexcerpt}
A certified spread-only report gives the known \CtwoSD{population} \CtwoSD{standard} \CtwoSD{deviation} of delivery fees as \CtwoSD{8} \CtwoSD{thousand} \CtwoSD{COP}, distinct from either \SSD{sample} \SSD{standard} \SSD{deviation}; the \CtwoMean{population} \CtwoMean{mean} is initially withheld. The initial study uses \CtwoConfidence{95\%} \CtwoConfidence{confidence} and the supplied \CtwoCritical{positive} \CtwoCritical{standard} \CtwoCritical{normal} \CtwoCritical{critical} \CtwoCritical{value} \CtwoCritical{1.960}.
\end{probleminfoexcerpt}

\begin{probleminfoexcerpt}
Assume the fee distribution is sufficiently regular with no dominant outliers, so the normal approximation to sample means is suitable. Each sample is selected without replacement from the full register. Study selections are independent and may overlap. Each study samples 0.064\% of the population, so the finite-population correction is negligible.
\end{probleminfoexcerpt}

\textbf{Later investigation (unavailable for initial estimation).}

\begin{probleminfoexcerpt}
A hypothetical later comprehensive investigation of the entire same-day delivery population establishes the \CtwoMean{population} \CtwoMean{mean} as \CtwoMean{33} \CtwoMean{thousand} \CtwoMean{COP}; this is not established by the repeated samples.
\end{probleminfoexcerpt}

The shared \CtwoPoint{point estimate} is $\CtwoPoint{\bar x}=\CtwoPoint{32}$ thousand COP, the average delivery fee in the initial sample; it estimates the population mean with a single value.

\textbf{\CtwoConfidence{95\%} confidence}

\keytable{p{.10\linewidth}p{.17\linewidth}p{.32\linewidth}K}{
\textbf{Symbol}&\textbf{Concept}&\textbf{Calculation}&\textbf{Result and meaning in context}\\\midrule
$\CtwoPoint{\bar x}$&\CtwoPoint{Point estimate}&$\CtwoPoint{\bar x}$=$\CtwoPoint{32}$&$\CtwoPoint{32}$ thousand COP; estimated population mean delivery fee.\\
$\CtwoMoE{\mathrm{MoE}}$&\CtwoMoE{Margin of error}&$\CtwoCritical{z_{1-\alpha/2}}\frac{\CtwoSD{\sigma}}{\sqrt{\CtwoSize{n}}}=\CtwoCritical{1.960}\frac{\CtwoSD{8}}{\sqrt{\CtwoSize{64}}}$&$\CtwoMoE{1.96}$ thousand COP; distance from estimate to either endpoint.\\
$\CtwoEndpoint{L}$&\CtwoEndpoint{Lower endpoint}&$\CtwoPoint{\bar x}-\CtwoMoE{\mathrm{MoE}}=\CtwoPoint{32}-\CtwoMoE{1.96}$&$\CtwoEndpoint{30.04}$ thousand COP; lower bound for the population mean.\\
$\CtwoEndpoint{U}$&\CtwoEndpoint{Upper endpoint}&$\CtwoPoint{\bar x}+\CtwoMoE{\mathrm{MoE}}=\CtwoPoint{32}+\CtwoMoE{1.96}$&$\CtwoEndpoint{33.96}$ thousand COP; upper bound for the population mean.\\
$\CtwoCI{\mathrm{CI}}$&\CtwoCI{Confidence interval}&$[\CtwoEndpoint{L},\CtwoEndpoint{U}]=\CtwoCI{[30.04,33.96]}$&Plausible population mean delivery fee values, in thousand COP.\\
$\CtwoWidth{W}$&\CtwoWidth{Full width}&$\CtwoEndpoint{U}-\CtwoEndpoint{L}=\CtwoEndpoint{33.96}-\CtwoEndpoint{30.04}$
$=2(\CtwoMoE{1.96})=\CtwoWidth{3.92}$&$\CtwoWidth{3.92}$ thousand COP; smaller width means greater precision.\\}



We are \CtwoConfidence{95\%} confident that the mean delivery fee for all 100,000 completed deliveries by registered couriers in a Bogota platform on November 2, 2026 lies between $\CtwoEndpoint{30.04}$ and $\CtwoEndpoint{33.96}$ thousand COP, in the repeated-sampling sense described in Task 3. This estimates a population mean, not a range containing \CtwoConfidence{95\%} of individual values.

\secInicio[][sec:p3]

\clearpage

\taskprompt{C2 Task 3}{2}{\cTwoC}

\worked{C2 Task 3 \score{2}}{sol:p3c2c}

\textbf{Relevant problem information}

\begin{probleminfoexcerpt}
A platform investigates all \Pop{100,000} \Pop{completed} \Pop{deliveries} by its registered couriers in Bogota on November 2, 2026. The measurement is the fee for each delivery in thousand COP.
\end{probleminfoexcerpt}

\begin{probleminfoexcerpt}
A certified spread-only report gives the known \CtwoSD{population} \CtwoSD{standard} \CtwoSD{deviation} of delivery fees as \CtwoSD{8} \CtwoSD{thousand} \CtwoSD{COP}, distinct from either \SSD{sample} \SSD{standard} \SSD{deviation}; the \CtwoMean{population} \CtwoMean{mean} is initially withheld.
\end{probleminfoexcerpt}

\begin{probleminfoexcerpt}
Assume the fee distribution is sufficiently regular with no dominant outliers, so the normal approximation to sample means is suitable. Each sample is selected without replacement from the full register. Study selections are independent and may overlap. Each study samples 0.064\% of the population, so the finite-population correction is negligible.
\end{probleminfoexcerpt}

\begin{probleminfoexcerpt}
Five subsequent independent studies, separate from the two original samples, use \CtwoSize{64} deliveries each, the same known \CtwoSD{population} \CtwoSD{standard} \CtwoSD{deviation} and the same \CtwoConfidence{95\%} normal procedure.
\end{probleminfoexcerpt}

\begin{probleminfoexcerpt}
A hypothetical later comprehensive investigation of the entire same-day delivery population establishes the \CtwoMean{population} \CtwoMean{mean} as \CtwoMean{33} \CtwoMean{thousand} \CtwoMean{COP}; this is not established by the repeated samples. The subsequent interval endpoints below are in thousand COP.
\end{probleminfoexcerpt}

\keytable{KKK}{
\textbf{Subsequent study}&\textbf{\CtwoEndpoint{Lower endpoint}}&\textbf{\CtwoEndpoint{Upper endpoint}}\\\midrule
1&\CtwoEndpoint{30.04}&\CtwoEndpoint{33.96}\\
2&\CtwoEndpoint{31.04}&\CtwoEndpoint{34.96}\\
3&\CtwoEndpoint{32.04}&\CtwoEndpoint{35.96}\\
4&\CtwoEndpoint{33.04}&\CtwoEndpoint{36.96}\\
5&\CtwoEndpoint{30.54}&\CtwoEndpoint{34.46}\\
}

Test each subsequent interval against $\CtwoMean{\mu_{\mathrm{known}}}=\CtwoMean{33}$ thousand COP; include the mean when it equals an endpoint. All endpoints below use thousand COP.

\keytable{KKKK}{
\textbf{Study}&\textbf{Lower endpoint}&\textbf{Upper endpoint}&\textbf{Covers later mean?}\\\midrule
1&$\CtwoEndpoint{30.04}$&$\CtwoEndpoint{33.96}$&\CtwoCoverage{Yes}\\
2&$\CtwoEndpoint{31.04}$&$\CtwoEndpoint{34.96}$&\CtwoCoverage{Yes}\\
3&$\CtwoEndpoint{32.04}$&$\CtwoEndpoint{35.96}$&\CtwoCoverage{Yes}\\
4&$\CtwoEndpoint{33.04}$&$\CtwoEndpoint{36.96}$&\CtwoCoverage{No}\\
5&$\CtwoEndpoint{30.54}$&$\CtwoEndpoint{34.46}$&\CtwoCoverage{Yes}\\
}


At \CtwoConfidence{95\%}: $\CtwoCoverage{N_{\mathrm{cover}}}=4$ and $N(95\%)=5$ subsequent intervals. Thus
\[\begin{aligned}
\CtwoCoverage{\mathrm{Cov}}(\CtwoConfidence{95\%})&=\frac{\CtwoCoverage{4}}{\CtwoCoverage{5}}\,100\%=\CtwoCoverage{80\%},\\
\CtwoError{E_{\mathrm{abs}}}&=|\CtwoCoverage{80}-\CtwoConfidence{95}|=\CtwoError{15}\text{ percentage points}.
\end{aligned}\]

The \CtwoCoverage{observed coverage} of the later-established mean delivery fee is \CtwoError{15} percentage points below nominal \CtwoConfidence{95\%}. Each listed interval has full width $2(\CtwoMoE{1.96})=\CtwoWidth{3.92}$ thousand COP, consistent with the stated $n$, known $\CtwoSD{\sigma}$ and critical value. Its center is its own subsequent sample mean, found as $(L_i+U_i)/2$, not the initial mean.

\Needspace{14\baselineskip}
\textbf{Interpretation.} Under the stated model, repeated independent use of each confidence procedure covers the fixed population mean in approximately its nominal percentage of intervals over the long run. A finite collection of intervals need not have exactly nominal coverage, and successive observed proportions need not approach it monotonically. The discrepancies here can arise from finite-study variation; they do not by themselves invalidate the procedure or prove that nominal coverage is wrong.

The population mean is fixed. Once an interval is calculated, it either contains that mean or does not; no nominal probability is assigned to this already determined event. Coverage counts compare each subsequent study's own interval with the later-established population mean. They do not count sample means or individual observations inside the initial interval. The initial and comparison samples are excluded from the subsequent-study denominator.

\secInicio[][sec:p3]

\clearpage

\subsection*{Proposed partial-credit guidance: 12 points}
The task marks are prescribed by the exam; only the detailed component splits below are \textbf{proposed}. Assess the student's own coherent scenario; the model answer is not a required topic or dataset.
\keytable{Kp{.11\linewidth}K}{
\textbf{Task}&\textbf{Points}&\textbf{Analytic evidence, including construction validity}\\\midrule
C1 Task 1&3&0.75 complete fixed population scope and sizes; 0.75 two clearly defined independent samples with feasible summaries and valid binary counts; 0.5 mean/SD pairings, units and distinct population/sample spreads; 1 both proportion calculations and interpretations, plus correct known/unknown status.\\
C1 Task 2&3&1.5 three correct signed differences (0.5 each; percentage points for proportions); 1.5 fixed population characteristics and sound explanation of sample variation.\\
C2 Task 1&2&0.5 adequate assumptions, negligible fractions, same-population study design and known spread; 0.75 completed parameter/sample groups, initially unknown mean and explicitly later comprehensive investigation; 0.75 SE and complete confidence/significance/tails/critical-value settings.\\
C2 Task 2&2&0.4 point estimate; 0.4 correct margin from known SD; 0.6 endpoints and interval with substitutions; 0.3 full width; 0.3 units and contextual mean interpretation.\\
C2 Task 3&2&0.75 complete, valid subsequent evidence consistent with the procedure and later mean; 0.75 counts, coverage and absolute percentage-point difference; 0.5 long-run/finite-study explanation with no fixed-mean probability error.\\}
\textbf{Rubric reconciliation:} C1: $3+3=6$; C2: $2+2+2=6$; overall: $6+6=12$. No separate construction bonus is available. Missing or invalid scenario elements lose credit within the components that assess them. Correct conditional methods may earn credit even when one supplied datum is flawed; do not repeatedly penalize the same arithmetic mistake. Award proportional credit within components for partially correct work. Other internally consistent scenarios and correct solutions are acceptable.

\textbf{Whole-exam scoring check.} C1: $4+4+6=14$; C2: $4+4+6=14$; overall: $8+8+12=28$. Pass each criterion independently with at least 4 out of 14 points; scores never transfer.


\secInicio[][sec:p3]


\clearpage
\section{Summarized Solutions}
\label{sec:summary}
\subsection*{Problem 1: Retail electricity costs}
\label{summary:p1}

\Needspace{7\baselineskip}
\par\textbf{\hyperref[sol:p1c1a]{C1 Task 1}}\par
\Pop{Population}: \Pop{200,000 registered retail businesses} operating in Medellin throughout October 2026; measurement: October electricity bill; condition: electronic invoicing in October. Both \Sam{samples} are independent random selections from this same population. All monetary entries below are in \textbf{thousand COP}.
Here $\CtwoMean{\mu}$ is the \CtwoMean{mean October} \CtwoMean{electricity bill}; $\CtwoSD{\sigma}$ is the \CtwoSD{population standard} \CtwoSD{deviation}, the \CtwoSD{spread of October} \CtwoSD{electricity bills}; and $\PP{p}$ is the \PP{proportion of all} \PP{eligible businesses} \PP{that used electronic} \PP{invoicing for their} \PP{October electricity} \PP{bills}.
\keytable{p{.29\linewidth}KK}{
\textbf{Parameter--statistic pair} & \textbf{Population status} & \textbf{Known initial A; second B}\\\midrule
\Pop{Size} / \Sam{sizes} & $\Pop{N}$ = $\Pop{200,000}$ (known) & $\Sam{n_A}=\Sam{n_B}$ = $\Sam{100}$\\
\CtwoMean{Mean} / \CtwoPoint{sample means} & $\CtwoMean{\mu}$: initially unknown & $\CtwoPoint{\bar x_A}$ = $\CtwoPoint{240}$; $\CtwoPoint{\bar x_B}$ = $\CtwoPoint{252}$\\
\CtwoSD{SD} / \SSD{sample SDs} & $\CtwoSD{\sigma}$ = $\CtwoSD{60}$ (known) & $\SSD{s_A}$ = $\SSD{58}$; $\SSD{s_B}$ = $\SSD{66}$\\
\PP{Proportion} / \SP{sample proportions} & $\PP{p}$: unknown & $\SP{\hat p_A}$ = $\Sam{40}$/$\Sam{100}$ = $\SP{0.4}=\SP{40\%}$\newline $\SP{\hat p_B}$ = $\Sam{55}$/$\Sam{100}$ = $\SP{0.55}=\SP{55\%}$\\}

\Needspace{7\baselineskip}
\par\textbf{\hyperref[sol:p1c1b]{C1 Task 2}}\par
\keytable{p{.28\linewidth}Kp{.26\linewidth}}{
\textbf{Quantity} & \textbf{Second-minus-initial calculation} & \textbf{Signed change and units}\\\midrule
\CtwoPoint{Sample mean} $\CtwoPoint{\Delta\bar x}$ & $\CtwoPoint{\bar x_B}-\CtwoPoint{\bar x_A}=\CtwoPoint{252}-\CtwoPoint{240}$ & $\CtwoPoint{+12}$ thousand COP\\
\SSD{Sample standard} \SSD{deviation} $\SSD{\Delta s}$ & $\SSD{s_B}-\SSD{s_A}=\SSD{66}-\SSD{58}$ & $\SSD{+8}$ thousand COP\\
\SP{Sample proportion} $\SP{\Delta\hat p}$ & $100(\SP{\hat p_B}-\SP{\hat p_A})$\newline $=100(\SP{0.55}-\SP{0.4})$ & $\SP{+15}$ percentage points\\}
Fixed population scope keeps $\CtwoMean{\mu}$, $\CtwoSD{\sigma}$ and $\PP{p}$ fixed; differences between sample statistics reflect sampling variation, not established parameter changes.

\Needspace{22\baselineskip}
\par\textbf{\hyperref[sol:p1c2a]{C2 Task 1}}\par
\keytable{p{.32\linewidth}p{.14\linewidth}K}{
\textbf{Quantity} & \textbf{Notation} & \textbf{Value or status}\\\midrule
\multicolumn{3}{l}{\textbf{Population parameters --- initial information}}\\
\CtwoMean{Population mean} & $\CtwoMean{\mu}$ & Initially unknown \CtwoMean{mean October} \CtwoMean{electricity bill}; the target of estimation.\\
\CtwoSD{Population standard} \CtwoSD{deviation} & $\CtwoSD{\sigma}$ & $\CtwoSD{60}$ thousand COP (known); \CtwoSD{spread of October} \CtwoSD{electricity bills}.\\\midrule
\multicolumn{3}{l}{\textbf{Population parameter --- later investigation only}}\\
\CtwoMean{Later-established mean} & $\CtwoMean{\mu_{\mathrm{known}}}$ & $\CtwoMean{245}$ thousand COP; the same fixed \CtwoMean{mean October} \CtwoMean{electricity bill}, $\CtwoMean{\mu}$, established by the comprehensive population investigation. Unavailable for initial estimation.\\\midrule
\multicolumn{3}{l}{\textbf{Initial sample inputs --- sample A only}}\\
\CtwoSize{Sample size} & $\CtwoSize{n}$ & $\CtwoSize{100}$ businesses.\\
\CtwoPoint{Sample mean} & $\CtwoPoint{\bar x}$ & $\CtwoPoint{240}$ thousand COP.\\
\SSD{Sample standard} \SSD{deviation} & $\SSD{s}$ & $\SSD{58}$ thousand COP.\\\midrule
\CtwoSE{Standard error} & $\CtwoSE{\mathrm{SE}}$ & $\CtwoSE{\mathrm{SE}}=\dfrac{\CtwoSD{\sigma}}{\sqrt{\CtwoSize{n}}}=\dfrac{\CtwoSD{60}}{\sqrt{\CtwoSize{100}}}=\CtwoSE{6}$ thousand COP.\\}
\Needspace{7\baselineskip}
\textbf{Procedure settings}
\keytable{KKKK}{
\CtwoConfidence{Confidence} $\CtwoConfidence{C}$ & \CtwoSignificance{Significance} $\CtwoSignificance{\alpha}$ & \CtwoSignificance{Each tail} $\CtwoSignificance{\alpha/2}$ & \CtwoCritical{Critical value} $\CtwoCritical{z_{1-\alpha/2}}$\\\midrule
$\CtwoConfidence{95\%}$ & $\CtwoSignificance{0.05}$ & $\CtwoSignificance{0.025}$ & $\CtwoCritical{1.960}$\\}

\Needspace{12\baselineskip}
\par\textbf{\hyperref[sol:p1c2b]{C2 Task 2}}\par
\keytable{p{.30\linewidth}K}{
\textbf{Result (thousand COP)} & $\CtwoConfidence{95\%}$\\
\CtwoPoint{Point estimate} $\CtwoPoint{\bar x}$ & $\CtwoPoint{240}$\\
\CtwoMoE{Margin of error} $\CtwoMoE{\mathrm{MoE}}$ & $\CtwoMoE{11.76}$\\
\CtwoEndpoint{Lower endpoint} $\CtwoEndpoint{L}$ & $\CtwoEndpoint{228.24}$\\
\CtwoEndpoint{Upper endpoint} $\CtwoEndpoint{U}$ & $\CtwoEndpoint{251.76}$\\
\CtwoCI{Interval} $\CtwoCI{\mathrm{CI}}$ & $\boxed{\CtwoCI{[228.24,251.76]}}$\\
\CtwoWidth{Full width} $\CtwoWidth{W}$ & $\CtwoWidth{23.52}$\\}

\Needspace{7\baselineskip}
\par\textbf{\hyperref[sol:p1c2c]{C2 Task 3}}\par
\keytable{KKKK}{
\CtwoConfidence{Confidence} & \CtwoCoverage{Covering / total} & \CtwoCoverage{Observed coverage} & \CtwoError{Absolute difference} (pp)\\\midrule
$\CtwoConfidence{95\%}$ & $\CtwoCoverage{7}/\CtwoCoverage{10}$ & $\CtwoCoverage{70\%}$ & $\CtwoError{25}$\\}
Count subsequent intervals with $\CtwoEndpoint{L_i}\le \CtwoMean{\mu_{\mathrm{known}}}\le \CtwoEndpoint{U_i}$; exclude A and B.
\secInicio[][sec:p1]

\subsection*{Problem 2: Savings-account interest}
\label{summary:p2}

\Needspace{7\baselineskip}
\par\textbf{\hyperref[sol:p2c1a]{C1 Task 1}}\par
\Pop{Population}: \Pop{500,000 positive-balance personal savings accounts} in the Cartagena branch network on October 31, 2026; measurement: October interest credit; condition: mobile banking activated on that date. Both \Sam{samples} are independent random selections from this same population. All monetary entries below are in \textbf{thousand COP}.
Here $\CtwoMean{\mu}$ is the \CtwoMean{mean October} \CtwoMean{interest credit}; $\CtwoSD{\sigma}$ is the \CtwoSD{population standard} \CtwoSD{deviation}, the \CtwoSD{spread of October} \CtwoSD{interest credits}; and $\PP{p}$ is the \PP{proportion of all} \PP{eligible accounts} \PP{that had mobile banking} \PP{activated on October 31}.
\keytable{p{.29\linewidth}KK}{
\textbf{Parameter--statistic pair} & \textbf{Population status} & \textbf{Known initial A; second B}\\\midrule
\Pop{Size} / \Sam{sizes} & $\Pop{N}$ = $\Pop{500,000}$ (known) & $\Sam{n_A}=\Sam{n_B}$ = $\Sam{144}$\\
\CtwoMean{Mean} / \CtwoPoint{sample means} & $\CtwoMean{\mu}$: initially unknown & $\CtwoPoint{\bar x_A}$ = $\CtwoPoint{82}$; $\CtwoPoint{\bar x_B}$ = $\CtwoPoint{79}$\\
\CtwoSD{SD} / \SSD{sample SDs} & $\CtwoSD{\sigma}$ = $\CtwoSD{36}$ (known) & $\SSD{s_A}$ = $\SSD{33}$; $\SSD{s_B}$ = $\SSD{39}$\\
\PP{Proportion} / \SP{sample proportions} & $\PP{p}$: unknown & $\SP{\hat p_A}$ = $\Sam{54}$/$\Sam{144}$ = $\SP{0.375}=\SP{37.5\%}$\newline $\SP{\hat p_B}$ = $\Sam{72}$/$\Sam{144}$ = $\SP{0.5}=\SP{50\%}$\\}

\Needspace{7\baselineskip}
\par\textbf{\hyperref[sol:p2c1b]{C1 Task 2}}\par
\keytable{p{.28\linewidth}Kp{.26\linewidth}}{
\textbf{Quantity} & \textbf{Second-minus-initial calculation} & \textbf{Signed change and units}\\\midrule
\CtwoPoint{Sample mean} $\CtwoPoint{\Delta\bar x}$ & $\CtwoPoint{\bar x_B}-\CtwoPoint{\bar x_A}=\CtwoPoint{79}-\CtwoPoint{82}$ & $\CtwoPoint{-3}$ thousand COP\\
\SSD{Sample standard} \SSD{deviation} $\SSD{\Delta s}$ & $\SSD{s_B}-\SSD{s_A}=\SSD{39}-\SSD{33}$ & $\SSD{+6}$ thousand COP\\
\SP{Sample proportion} $\SP{\Delta\hat p}$ & $100(\SP{\hat p_B}-\SP{\hat p_A})$\newline $=100(\SP{0.5}-\SP{0.375})$ & $\SP{+12.5}$ percentage points\\}
Fixed population scope keeps $\CtwoMean{\mu}$, $\CtwoSD{\sigma}$ and $\PP{p}$ fixed; differences between sample statistics reflect sampling variation, not established parameter changes.


\Needspace{22\baselineskip}
\par\textbf{\hyperref[sol:p2c2a]{C2 Task 1}}\par
\keytable{p{.32\linewidth}p{.14\linewidth}K}{
\textbf{Quantity} & \textbf{Notation} & \textbf{Value or status}\\\midrule
\multicolumn{3}{l}{\textbf{Population parameters --- initial information}}\\
\CtwoMean{Population mean} & $\CtwoMean{\mu}$ & Initially unknown \CtwoMean{mean October} \CtwoMean{interest credit}; the target of estimation.\\
\CtwoSD{Population standard} \CtwoSD{deviation} & $\CtwoSD{\sigma}$ & $\CtwoSD{36}$ thousand COP (known); \CtwoSD{spread of October} \CtwoSD{interest credits}.\\\midrule
\multicolumn{3}{l}{\textbf{Population parameter --- later investigation only}}\\
\CtwoMean{Later-established mean} & $\CtwoMean{\mu_{\mathrm{known}}}$ & $\CtwoMean{84}$ thousand COP; the same fixed \CtwoMean{mean October} \CtwoMean{interest credit}, $\CtwoMean{\mu}$, established by the comprehensive population investigation. Unavailable for initial estimation.\\\midrule
\multicolumn{3}{l}{\textbf{Initial sample inputs --- sample A only}}\\
\CtwoSize{Sample size} & $\CtwoSize{n}$ & $\CtwoSize{144}$ accounts.\\
\CtwoPoint{Sample mean} & $\CtwoPoint{\bar x}$ & $\CtwoPoint{82}$ thousand COP.\\
\SSD{Sample standard} \SSD{deviation} & $\SSD{s}$ & $\SSD{33}$ thousand COP.\\\midrule
\CtwoSE{Standard error} & $\CtwoSE{\mathrm{SE}}$ & $\CtwoSE{\mathrm{SE}}=\dfrac{\CtwoSD{\sigma}}{\sqrt{\CtwoSize{n}}}=\dfrac{\CtwoSD{36}}{\sqrt{\CtwoSize{144}}}=\CtwoSE{3}$ thousand COP.\\}
\Needspace{7\baselineskip}
\textbf{Procedure settings}
\keytable{KKKK}{
\CtwoConfidence{Confidence} $\CtwoConfidence{C}$ & \CtwoSignificance{Significance} $\CtwoSignificance{\alpha}$ & \CtwoSignificance{Each tail} $\CtwoSignificance{\alpha/2}$ & \CtwoCritical{Critical value} $\CtwoCritical{z_{1-\alpha/2}}$\\\midrule
$\CtwoConfidence{90\%}$ & $\CtwoSignificance{0.10}$ & $\CtwoSignificance{0.05}$ & $\CtwoCritical{1.645}$\\
$\CtwoConfidence{99\%}$ & $\CtwoSignificance{0.01}$ & $\CtwoSignificance{0.005}$ & $\CtwoCritical{2.576}$\\}

\Needspace{12\baselineskip}
\par\textbf{\hyperref[sol:p2c2b]{C2 Task 2}}\par
\keytable{p{.30\linewidth}KK}{
\textbf{Result (thousand COP)} & $\CtwoConfidence{90\%}$ & $\CtwoConfidence{99\%}$\\
\CtwoPoint{Point estimate} $\CtwoPoint{\bar x}$ & $\CtwoPoint{82}$ & $\CtwoPoint{82}$\\
\CtwoMoE{Margin of error} $\CtwoMoE{\mathrm{MoE}}$ & $\CtwoMoE{4.935}$ & $\CtwoMoE{7.728}$\\
\CtwoEndpoint{Lower endpoint} $\CtwoEndpoint{L}$ & $\CtwoEndpoint{77.065}$ & $\CtwoEndpoint{74.272}$\\
\CtwoEndpoint{Upper endpoint} $\CtwoEndpoint{U}$ & $\CtwoEndpoint{86.935}$ & $\CtwoEndpoint{89.728}$\\
\CtwoCI{Interval} $\CtwoCI{\mathrm{CI}}$ & $\boxed{\CtwoCI{[77.065,86.935]}}$ & $\boxed{\CtwoCI{[74.272,89.728]}}$\\
\CtwoWidth{Full width} $\CtwoWidth{W}$ & $\CtwoWidth{9.87}$ & $\CtwoWidth{15.456}$\\}

\Needspace{7\baselineskip}
\par\textbf{\hyperref[sol:p2c2c]{C2 Task 3}}\par
\keytable{KKKK}{
\CtwoConfidence{Confidence} & \CtwoCoverage{Covering / total} & \CtwoCoverage{Observed coverage} & \CtwoError{Absolute difference} (pp)\\\midrule
$\CtwoConfidence{90\%}$ & $\CtwoCoverage{176}/\CtwoCoverage{200}$ & $\CtwoCoverage{88\%}$ & $\CtwoError{2}$\\
$\CtwoConfidence{99\%}$ & $\CtwoCoverage{196}/\CtwoCoverage{200}$ & $\CtwoCoverage{98\%}$ & $\CtwoError{1}$\\}
At $\CtwoConfidence{99\%}$, covering count $\CtwoCoverage{176}+\CtwoCoverage{20}=\CtwoCoverage{196}$; the wider interval contains the $\CtwoConfidence{90\%}$ interval.
Count subsequent intervals with $\CtwoEndpoint{L_i}\le \CtwoMean{\mu_{\mathrm{known}}}\le \CtwoEndpoint{U_i}$; exclude A and B.
\secInicio[][sec:p2]

\Needspace{10\baselineskip}
\subsection*{Problem 3: Create an integrated scenario from scratch}
\label{summary:p3}

\Needspace{7\baselineskip}
\par\textbf{\hyperref[sol:p3c1a]{C1 Task 1}}\par
Existing hypothetical model: \Pop{100,000 completed deliveries} by registered platform couriers in Bogota on November 2, 2026; measurement: delivery fee; condition: optional parcel insurance. Both \Sam{samples} are independent random selections from this same population. All monetary entries below are in \textbf{thousand COP}.
Here $\CtwoMean{\mu}$ is the \CtwoMean{mean delivery fee}; $\CtwoSD{\sigma}$ is the \CtwoSD{population standard} \CtwoSD{deviation}, the \CtwoSD{spread of delivery fees}; and $\PP{p}$ is the \PP{proportion of all} \PP{eligible deliveries} \PP{that included optional} \PP{parcel insurance}.
\keytable{p{.29\linewidth}KK}{
\textbf{Parameter--statistic pair} & \textbf{Population status} & \textbf{Known initial A; second B}\\\midrule
\Pop{Size} / \Sam{sizes} & $\Pop{N}$ = $\Pop{100,000}$ (known) & $\Sam{n_A}=\Sam{n_B}$ = $\Sam{64}$\\
\CtwoMean{Mean} / \CtwoPoint{sample means} & $\CtwoMean{\mu}$: initially unknown & $\CtwoPoint{\bar x_A}$ = $\CtwoPoint{32}$; $\CtwoPoint{\bar x_B}$ = $\CtwoPoint{34}$\\
\CtwoSD{SD} / \SSD{sample SDs} & $\CtwoSD{\sigma}$ = $\CtwoSD{8}$ (known) & $\SSD{s_A}$ = $\SSD{7}$; $\SSD{s_B}$ = $\SSD{9}$\\
\PP{Proportion} / \SP{sample proportions} & $\PP{p}$: unknown & $\SP{\hat p_A}$ = $\Sam{24}$/$\Sam{64}$ = $\SP{0.375}=\SP{37.5\%}$\newline $\SP{\hat p_B}$ = $\Sam{32}$/$\Sam{64}$ = $\SP{0.5}=\SP{50\%}$\\}

\Needspace{7\baselineskip}
\par\textbf{\hyperref[sol:p3c1b]{C1 Task 2}}\par
\keytable{p{.28\linewidth}Kp{.26\linewidth}}{
\textbf{Quantity} & \textbf{Second-minus-initial calculation} & \textbf{Signed change and units}\\\midrule
\CtwoPoint{Sample mean} $\CtwoPoint{\Delta\bar x}$ & $\CtwoPoint{\bar x_B}-\CtwoPoint{\bar x_A}=\CtwoPoint{34}-\CtwoPoint{32}$ & $\CtwoPoint{+2}$ thousand COP\\
\SSD{Sample standard} \SSD{deviation} $\SSD{\Delta s}$ & $\SSD{s_B}-\SSD{s_A}=\SSD{9}-\SSD{7}$ & $\SSD{+2}$ thousand COP\\
\SP{Sample proportion} $\SP{\Delta\hat p}$ & $100(\SP{\hat p_B}-\SP{\hat p_A})$\newline $=100(\SP{0.5}-\SP{0.375})$ & $\SP{+12.5}$ percentage points\\}
Fixed population scope keeps $\CtwoMean{\mu}$, $\CtwoSD{\sigma}$ and $\PP{p}$ fixed; differences between sample statistics reflect sampling variation, not established parameter changes.


\Needspace{22\baselineskip}
\par\textbf{\hyperref[sol:p3c2a]{C2 Task 1}}\par
\keytable{p{.32\linewidth}p{.14\linewidth}K}{
\textbf{Quantity} & \textbf{Notation} & \textbf{Value or status}\\\midrule
\multicolumn{3}{l}{\textbf{Population parameters --- initial information}}\\
\CtwoMean{Population mean} & $\CtwoMean{\mu}$ & Initially unknown \CtwoMean{mean delivery fee}; the target of estimation.\\
\CtwoSD{Population standard} \CtwoSD{deviation} & $\CtwoSD{\sigma}$ & $\CtwoSD{8}$ thousand COP (known); \CtwoSD{spread of delivery fees}.\\\midrule
\multicolumn{3}{l}{\textbf{Population parameter --- later investigation only}}\\
\CtwoMean{Later-established mean} & $\CtwoMean{\mu_{\mathrm{known}}}$ & $\CtwoMean{33}$ thousand COP; the same fixed \CtwoMean{mean delivery fee}, $\CtwoMean{\mu}$, established by the comprehensive population investigation. Unavailable for initial estimation.\\\midrule
\multicolumn{3}{l}{\textbf{Initial sample inputs --- sample A only}}\\
\CtwoSize{Sample size} & $\CtwoSize{n}$ & $\CtwoSize{64}$ deliveries.\\
\CtwoPoint{Sample mean} & $\CtwoPoint{\bar x}$ & $\CtwoPoint{32}$ thousand COP.\\
\SSD{Sample standard} \SSD{deviation} & $\SSD{s}$ & $\SSD{7}$ thousand COP.\\\midrule
\CtwoSE{Standard error} & $\CtwoSE{\mathrm{SE}}$ & $\CtwoSE{\mathrm{SE}}=\dfrac{\CtwoSD{\sigma}}{\sqrt{\CtwoSize{n}}}=\dfrac{\CtwoSD{8}}{\sqrt{\CtwoSize{64}}}=\CtwoSE{1}$ thousand COP.\\}
\Needspace{7\baselineskip}
\textbf{Procedure settings}
\keytable{KKKK}{
\CtwoConfidence{Confidence} $\CtwoConfidence{C}$ & \CtwoSignificance{Significance} $\CtwoSignificance{\alpha}$ & \CtwoSignificance{Each tail} $\CtwoSignificance{\alpha/2}$ & \CtwoCritical{Critical value} $\CtwoCritical{z_{1-\alpha/2}}$\\\midrule
$\CtwoConfidence{95\%}$ & $\CtwoSignificance{0.05}$ & $\CtwoSignificance{0.025}$ & $\CtwoCritical{1.960}$\\}

\Needspace{12\baselineskip}
\par\textbf{\hyperref[sol:p3c2b]{C2 Task 2}}\par
\keytable{p{.30\linewidth}K}{
\textbf{Result (thousand COP)} & $\CtwoConfidence{95\%}$\\
\CtwoPoint{Point estimate} $\CtwoPoint{\bar x}$ & $\CtwoPoint{32}$\\
\CtwoMoE{Margin of error} $\CtwoMoE{\mathrm{MoE}}$ & $\CtwoMoE{1.96}$\\
\CtwoEndpoint{Lower endpoint} $\CtwoEndpoint{L}$ & $\CtwoEndpoint{30.04}$\\
\CtwoEndpoint{Upper endpoint} $\CtwoEndpoint{U}$ & $\CtwoEndpoint{33.96}$\\
\CtwoCI{Interval} $\CtwoCI{\mathrm{CI}}$ & $\boxed{\CtwoCI{[30.04,33.96]}}$\\
\CtwoWidth{Full width} $\CtwoWidth{W}$ & $\CtwoWidth{3.92}$\\}

\Needspace{7\baselineskip}
\par\textbf{\hyperref[sol:p3c2c]{C2 Task 3}}\par
\keytable{KKKK}{
\CtwoConfidence{Confidence} & \CtwoCoverage{Covering / total} & \CtwoCoverage{Observed coverage} & \CtwoError{Absolute difference} (pp)\\\midrule
$\CtwoConfidence{95\%}$ & $\CtwoCoverage{4}/\CtwoCoverage{5}$ & $\CtwoCoverage{80\%}$ & $\CtwoError{15}$\\}
Count subsequent intervals with $\CtwoEndpoint{L_i}\le \CtwoMean{\mu_{\mathrm{known}}}\le \CtwoEndpoint{U_i}$; exclude A and B.
Long-run coverage is approximately nominal confidence; finite-study variation can explain these differences. The fixed \CtwoMean{population mean} either lies in a calculated interval or does not: confidence is not its probability, nor a count of sample means inside the initial interval.
\secInicio[][sec:p3]

\end{document}
